\documentclass[12pt,a4paper]{article}

%% ── Packages ────────────────────────────────────────────────────────────────
\usepackage[T1]{fontenc}
\usepackage[utf8]{inputenc}
\usepackage[ngerman]{babel}

\usepackage{geometry}
\geometry{a4paper, margin=2.5cm}

\usepackage{amsmath,amssymb,amsthm}
\usepackage{mathtools}
\usepackage{thmtools}

\usepackage{algorithm}
\usepackage{algpseudocode}
\usepackage{listings}
\usepackage{xcolor}

\usepackage{tikz}
\usetikzlibrary{arrows.meta, positioning, shapes.geometric, calc, backgrounds, fit}

\usepackage{booktabs}
\usepackage{tabularx}
\usepackage{multirow}
\usepackage{longtable}
\usepackage{array}

\usepackage{graphicx}
\usepackage{caption}
\usepackage{subcaption}
\usepackage{float}

\usepackage{enumitem}
\usepackage{parskip}
\setlength{\parskip}{6pt}

\usepackage[colorlinks=true, linkcolor=blue!70!black, citecolor=green!50!black,
urlcolor=blue!60!black]{hyperref}
\usepackage{cleveref}

%% ── Theorem environments ────────────────────────────────────────────────────
\theoremstyle{plain}
\newtheorem{theorem}{Theorem}[section]
\newtheorem{lemma}[theorem]{Lemma}
\newtheorem{corollary}[theorem]{Corollary}
\newtheorem{proposition}[theorem]{Proposition}

\theoremstyle{definition}
\newtheorem{definition}[theorem]{Definition}
\newtheorem{example}[theorem]{Example}

\theoremstyle{remark}
\newtheorem{remark}[theorem]{Remark}

%% ── Algorithm style ─────────────────────────────────────────────────────────
\algrenewcommand\algorithmicrequire{\textbf{Input:}}
\algrenewcommand\algorithmicensure{\textbf{Output:}}
\algnewcommand{\LineComment}[1]{\State \(\triangleright\) \textit{#1}}

%% ── Colour scheme ───────────────────────────────────────────────────────────
\definecolor{algoblue}{RGB}{30,90,160}
\definecolor{algogray}{RGB}{245,245,245}
\definecolor{commentgreen}{RGB}{0,120,50}
\definecolor{keywordred}{RGB}{160,30,30}

%% ── Metadata ────────────────────────────────────────────────────────────────
\title{\textbf{The Hierarchical Subgraph-DFS Algorithm}\\[6pt]
	\large Formal Design, Complexity Analysis, and\\
	Cross-Domain Application Study}
\author{Stephan Epp}
\date{\today}

%% ════════════════════════════════════════════════════════════════════════════
\begin{document}
	\maketitle
	\tableofcontents
	\newpage
	
	%% ════════════════════════════════════════════════════════════════════════════
	\section{Introduction}
	\label{sec:intro}
	
	Modern computational problems across science, industry, and engineering share
	a fundamental structural property: the systems of interest are not flat, but
	\emph{hierarchically composed} -- they are systems of systems, networks of
	networks, or graphs of graphs.  The human connectome is a graph of neural
	sub-circuits; a gene-regulatory network is a graph of pathway modules; a
	logistics network is a graph of regional sub-networks; an industrial
	automation system is a graph of interacting Petri-net models.
	
	The central question that unifies these domains is:
	\begin{quote}
		\emph{Given a collection $\mathcal{G} = \{G_1, G_2, \ldots, G_N\}$ of
			graphs, which members stand in a subgraph relationship to which other
			members, and how does structural reachability propagate through that
			relationship?}
	\end{quote}
	
	This paper introduces the \textbf{Hierarchical Subgraph-DFS (HS-DFS) algorithm},
	a two-phase procedure that (i) constructs a \emph{meta-graph} $\mathcal{M}$
	whose nodes are the individual graphs and whose directed edges encode detected
	subgraph containment, and (ii) executes a Depth-First Search on $\mathcal{M}$
	to generate a \emph{DFS forest} that exposes the full hierarchical structure of
	the collection.  The pairwise subgraph detection in Phase~(i) is performed
	by the \textbf{Subgraph Algorithm} \cite{epp2026subgraph}: a novel
	signature-based procedure that represents graphs as adjacency matrices,
	encodes each column as a unique integer via a polynomial hashing scheme,
	and decides containment by testing $n$ cyclic rotations of the
	column-signature sequence using Longest Common Subsequence (LCS)
	comparison -- achieving $\mathcal{O}(n^3)$ time without any exponential
	permutation search.
	
	The choice of DFS in Phase~2 is deliberate and algorithmically justified:
	unlike Breadth-First Search, DFS yields a \emph{DFS forest} -- a spanning
	forest that classifies every edge as a tree, back, forward, or cross edge,
	thereby revealing cycles, hierarchical depth, and strong-connectivity
	structure that BFS cannot expose \cite{cormen2009}.
	
	\paragraph{Contributions.} The main contributions of this paper are:
	\begin{enumerate}[label=(\roman*)]
		\item A complete formal specification of the \textbf{Subgraph Algorithm}
		\cite{epp2026subgraph} -- a signature-based, cyclic-rotation oracle for
		pairwise subgraph detection running in $\mathcal{O}(n^3)$, serving as
		the concrete oracle component of HS-DFS.
		\item A rigorous formal definition of the HS-DFS algorithm with full
		pseudocode, integrating the Subgraph Algorithm as its oracle.
		\item A comprehensive worst-case and average-case complexity analysis
		establishing $\mathcal{O}(n^5)$ time complexity for dense meta-graphs.
		\item A systematic survey of application domains illustrating how HS-DFS
		provides a unifying algorithmic primitive across neuroscience, genomics,
		logistics, industrial automation, software engineering, and beyond.
	\end{enumerate}
	
	%% ════════════════════════════════════════════════════════════════════════════
	\section{Preliminaries and Formal Definitions}
	\label{sec:prelim}
	
	\begin{definition}[Graph]
		A \emph{graph} $G = (V, E)$ consists of a finite set of vertices $V$ and a
		set of edges $E \subseteq V \times V$.  We write $|G|$ for $|V|$.
	\end{definition}
	
	\begin{definition}[Subgraph]
		A graph $H = (V_H, E_H)$ is a \emph{subgraph} of $G = (V_G, E_G)$,
		written $H \subseteq G$, if and only if $V_H \subseteq V_G$ and
		$E_H \subseteq E_G$.  We additionally call $H$ an \emph{induced subgraph}
		if $E_H = \{(u,v) \in E_G \mid u,v \in V_H\}$.
	\end{definition}
	
	\begin{definition}[Subgraph-Isomorphism]
		A graph $H = (V_H, E_H)$ is \emph{subgraph-isomorphic} to $G = (V_G, E_G)$,
		written $H \hookrightarrow G$, if there exists an injective mapping
		$\varphi: V_H \to V_G$ such that
		$(u,v) \in E_H \Rightarrow (\varphi(u), \varphi(v)) \in E_G$.
	\end{definition}
	
	\begin{definition}[Subgraph Oracle]
		\label{def:oracle}
		A \emph{subgraph oracle} $\mathcal{S}$ is an algorithm that, given two
		graphs $H$ and $G$ each of size at most $n$, decides whether
		$H \hookrightarrow G$.  In this paper, $\mathcal{S}$ is instantiated
		with the \textbf{Subgraph Algorithm} (\Cref{sec:subgraph}), which
		achieves $T_{\mathcal{S}}(n) = \mathcal{O}(n^3)$ via column signatures
		and cyclic rotations of adjacency matrices.
	\end{definition}
	
	\begin{remark}
		Subgraph isomorphism is NP-complete in general \cite{cook1971}.  The
		$\mathcal{O}(n^3)$ bound of the Subgraph Algorithm is attained by
		restricting vertex re-labellings to $n$ cyclic rotations instead of
		$n!$ arbitrary permutations.  This preserves sequential edge structure
		and is sufficient for the structural-similarity applications targeted
		by HS-DFS.  For fully general subgraph isomorphism (arbitrary label
		permutations), exponential running times remain necessary in the worst
		case.
	\end{remark}
	
	\begin{definition}[Graph Collection and Meta-Graph]
		Let $\mathcal{G} = \{G_1, \ldots, G_N\}$ be a collection of $N$ graphs.
		The \emph{meta-graph} $\mathcal{M} = (\mathcal{G},\, \mathcal{E})$ is the
		directed graph with node set $\mathcal{G}$ and edge set
		\[
		\mathcal{E} \;=\; \bigl\{(G_i, G_j) \;\mid\; G_i \hookrightarrow G_j,\;
		i \neq j\bigr\}.
		\]
	\end{definition}
	
	\begin{definition}[DFS Forest]
		Given a directed graph $\mathcal{M}$, a \emph{DFS forest} $\mathcal{F}$ is
		the spanning forest produced by Depth-First Search on $\mathcal{M}$.  Each
		edge in $\mathcal{M}$ is classified as a
		\emph{tree edge}, \emph{back edge}, \emph{forward edge}, or
		\emph{cross edge} with respect to $\mathcal{F}$.
	\end{definition}
	
	%% ════════════════════════════════════════════════════════════════════════════
	\section{The Subgraph Algorithm}
	\label{sec:subgraph}

	The HS-DFS framework requires a reliable, efficient pairwise subgraph oracle
	$\mathcal{S}$.  This section presents the \textbf{Subgraph Algorithm}
	\cite{epp2026subgraph} in full: a concrete procedure that decides
	$G \hookrightarrow G'$ by exploiting \emph{column signatures} of adjacency
	matrices and \emph{cyclic vertex re-labelling}, achieving
	$\mathcal{O}(n^3)$ time and $\mathcal{O}(n^2)$ space without recourse to
	exponential permutation enumeration.

	\subsection{Adjacency Matrix Representation}

	\begin{definition}[Adjacency Matrix]
		\label{def:adjmat}
		The \emph{adjacency matrix} of a graph $G = (V, E)$ with $n = |V|$ vertices
		is the matrix $A \in \{0,1\}^{n \times n}$ defined by
		\[
		A_{ij} = \begin{cases}
			1 & \text{if } (v_i, v_j) \in E, \\
			0 & \text{otherwise.}
		\end{cases}
		\]
	\end{definition}

	\begin{definition}[Subgraph-Decision Problem]
		\label{def:subgraph_decision}
		Given graphs $G$ and $G'$ with adjacency matrices
		$A \in \{0,1\}^{n_A \times n_A}$ and $B \in \{0,1\}^{n_B \times n_B}$,
		the \emph{subgraph-decision problem} asks whether $G \hookrightarrow G'$
		and determines the \emph{containment relation}:
		\begin{itemize}
			\item $B \supseteq A$: discard $G$; retain $G'$ as the more informative graph.
			\item $A \supseteq B$: discard $G'$; retain $G$.
			\item $A = B$ (isomorphic): retain either.
			\item Neither: retain both graphs as structurally incomparable.
		\end{itemize}
	\end{definition}

	\subsection{Column Signatures}

	\begin{definition}[Column Signature]
		\label{def:signature}
		Let $M \in \{0,1\}^{n \times n}$ be an adjacency matrix.  The
		\emph{column signature} of column $j \in \{0, \ldots, n-1\}$ is
		\[
		\sigma_j(M) \;=\; \sum_{i=0}^{n-1} M_{ij} \cdot 2^i \;+\; j \cdot 2^n.
		\]
		The \emph{signature sequence} of $M$ is
		$\boldsymbol{\sigma}(M) = [\sigma_0(M),\; \sigma_1(M),\; \ldots,\; \sigma_{n-1}(M)]$.
		The \emph{row-component sequence} strips the column-index weight:
		$r_j(M) = \sigma_j(M) \bmod 2^n$, i.e.,
		$r_j(M) = \sum_{i=0}^{n-1} M_{ij} \cdot 2^i$.
	\end{definition}

	\begin{lemma}[Injectivity of Column Signatures]
		\label{lem:injective}
		The column signature function $\sigma_j : \{0,1\}^{n \times n} \to \mathbb{N}$
		is injective: distinct column vectors at distinct positions always receive
		distinct signatures.
	\end{lemma}

	\begin{proof}
		The row component $\sum_{i=0}^{n-1} M_{ij} \cdot 2^i$ encodes the binary
		column vector uniquely as an integer in $[0, 2^n - 1]$.  The column-index
		weight $j \cdot 2^n$ maps each $j \in \{0, \ldots, n-1\}$ to the interval
		$[j \cdot 2^n,\; (j+1) \cdot 2^n - 1]$, which is disjoint from all other
		such intervals.  Hence two (column-vector, position) pairs produce the
		same signature only if both the binary vector and the column index coincide,
		establishing injectivity.
	\end{proof}

	\begin{example}
		\label{ex:signatures}
		Let $n = 4$.  For column $j = 0$ with column vector $[1, 0, 1, 0]^{\top}$:
		\[
		\sigma_0 = 1{\cdot}2^0 + 0{\cdot}2^1 + 1{\cdot}2^2 + 0{\cdot}2^3
		          + 0{\cdot}2^4 = 5, \qquad r_0 = 5.
		\]
		For column $j = 1$ with the identical vector $[1, 0, 1, 0]^{\top}$:
		\[
		\sigma_1 = 1{\cdot}2^0 + 0{\cdot}2^1 + 1{\cdot}2^2 + 0{\cdot}2^3
		          + 1{\cdot}2^4 = 21, \qquad r_1 = 5.
		\]
		The signatures differ ($5 \neq 21$) despite identical column vectors,
		confirming injectivity.  The equal row components ($r_0 = r_1 = 5$) reflect
		that both columns encode the same edge pattern -- used later to detect
		structural matches across different vertex positions.
	\end{example}

	\subsection{Cyclic Rotation Instead of Full Permutation}

	A cornerstone of the Subgraph Algorithm is the replacement of all $n!$
	vertex permutations by only $n$ cyclic rotations, reducing complexity from
	super-exponential to cubic while preserving structural correctness.

	\begin{definition}[Cyclic Rotation]
		\label{def:rotation}
		For a sequence $\mathbf{s} = [s_0, s_1, \ldots, s_{n-1}]$ and offset
		$k \in \{0, \ldots, n-1\}$, the \emph{$k$-th cyclic rotation} is
		\[
		\mathrm{rotate}_k(\mathbf{s})
		= [s_k,\; s_{k+1},\; \ldots,\; s_{n-1},\; s_0,\; \ldots,\; s_{k-1}].
		\]
	\end{definition}

	\begin{example}
		For $n = 4$ columns with row-component sequence $[r_0, r_1, r_2, r_3]$:
		\begin{align*}
			k = 0{:}\quad & [r_0,\; r_1,\; r_2,\; r_3] \\
			k = 1{:}\quad & [r_1,\; r_2,\; r_3,\; r_0] \\
			k = 2{:}\quad & [r_2,\; r_3,\; r_0,\; r_1] \\
			k = 3{:}\quad & [r_3,\; r_0,\; r_1,\; r_2]
		\end{align*}
		This yields only $4$ cases to examine rather than $4! = 24$ permutations.
	\end{example}

	\subsection{Subgraph Criterion via LCS}

	\begin{theorem}[Subgraph Criterion]
		\label{thm:subgraph_criterion}
		Let $G$ and $G'$ have row-component sequences
		$\mathbf{r}^A = [r_0^A, \ldots, r_{n_A-1}^A]$ and
		$\mathbf{r}^B = [r_0^B, \ldots, r_{n_B-1}^B]$.
		Then $G \hookrightarrow G'$ if and only if there exists
		$k \in \{0, \ldots, n_B - 1\}$ such that
		\[
		\mathrm{LCS}\!\bigl(\mathbf{r}^A,\;
		\mathrm{rotate}_k(\mathbf{r}^B)\bigr) \;\geq\; 2,
		\]
		where $\mathrm{LCS}$ denotes the length of the Longest Common Subsequence.
	\end{theorem}

	\begin{proof}
		A subgraph relationship $G \hookrightarrow G'$ requires at least two
		vertices connected by a common edge, i.e., at least two structurally
		matching column patterns must exist in the respective adjacency matrices.
		The row component $r_j(M)$ encodes precisely the edge pattern of column $j$
		without any position bias.  A cyclic rotation of $\mathbf{r}^B$ by offset
		$k$ models a cyclic re-labelling of the vertex indices of $G'$, which
		preserves all edge adjacencies while changing the column ordering.  By
		\Cref{lem:injective}, equal row components certify identical edge patterns.
		Therefore, an LCS of length $\geq 2$ between $\mathbf{r}^A$ and a rotation
		of $\mathbf{r}^B$ certifies that at least one edge of $G$ is present in
		$G'$ under the corresponding cyclic labelling, establishing
		$G \hookrightarrow G'$.
	\end{proof}

	\subsection{Pseudocode of the Subgraph Algorithm}

	\begin{algorithm}[H]
		\caption{Subgraph Algorithm$(A,\, B)$}
		\label{alg:subgraph}
		\begin{algorithmic}[1]
			\Require Adjacency matrices $A \in \{0,1\}^{n_A \times n_A}$,
			$B \in \{0,1\}^{n_B \times n_B}$
			\Ensure $\mathtt{true}$ iff $G_A \hookrightarrow G_B$

			\If{$n_A > n_B$}
			\Return $\mathtt{false}$
			\LineComment{Containment requires $n_A \leq n_B$}
			\EndIf

			\Statex
			\Statex \textbf{Step 1 -- Column signatures of $A$}
			\For{$j \leftarrow 0$ \textbf{to} $n_A - 1$}
			\State $\sigma_j^A \;\leftarrow\; \sum_{i=0}^{n_A-1} A_{ij} \cdot 2^i
			+ j \cdot 2^{n_A}$
			\EndFor

			\Statex
			\Statex \textbf{Step 2 -- Column signatures of $B$}
			\For{$j \leftarrow 0$ \textbf{to} $n_B - 1$}
			\State $\sigma_j^B \;\leftarrow\; \sum_{i=0}^{n_B-1} B_{ij} \cdot 2^i
			+ j \cdot 2^{n_B}$
			\EndFor

			\Statex
			\Statex \textbf{Step 3 -- Cyclic rotation match}
			\State $r_j^A \leftarrow \sigma_j^A \bmod 2^{n_A}$
			\quad for all $j \in \{0, \ldots, n_A - 1\}$
			\State $r_j^B \leftarrow \sigma_j^B \bmod 2^{n_B}$
			\quad for all $j \in \{0, \ldots, n_B - 1\}$

			\For{$k \leftarrow 0$ \textbf{to} $n_B - 1$}
			\State $\mathbf{s} \leftarrow
			\mathrm{rotate}_k\bigl([r_0^B,\; \ldots,\; r_{n_B-1}^B]\bigr)$
			\If{$\mathrm{LCS}\bigl([r_0^A, \ldots, r_{n_A-1}^A],\;
				\mathbf{s}\bigr) \geq 2$}
			\Return $\mathtt{true}$
			\EndIf
			\EndFor

			\Return $\mathtt{false}$
		\end{algorithmic}
	\end{algorithm}

	\subsection{Complexity and Correctness}

	\begin{theorem}[Time Complexity of the Subgraph Algorithm]
		\label{thm:subgraph_complexity}
		The Subgraph Algorithm runs in $\mathcal{O}(n^3)$ time, where
		$n = \max(n_A, n_B)$.
	\end{theorem}

	\begin{proof}
		\textbf{Steps~1 and~2:} Computing all $n_A$ (resp.\ $n_B$) column
		signatures requires iterating over a column vector of length $n_A$
		(resp.\ $n_B$) for each of the $n_A$ (resp.\ $n_B$) columns -- cost
		$\mathcal{O}(n^2)$ per step.

		\textbf{Step~3:} The outer loop executes at most $n_B \leq n$
		iterations.  For each rotation $k$, computing
		$\mathrm{rotate}_k$ costs $\mathcal{O}(n)$ and the LCS of two sequences
		of length at most $n$ is computed in $\mathcal{O}(n^2)$ by standard
		dynamic programming.  Total cost of Step~3:
		$\mathcal{O}(n \cdot (n + n^2)) = \mathcal{O}(n^3)$.

		\textbf{Total:}
		$\mathcal{O}(n^2) + \mathcal{O}(n^2) + \mathcal{O}(n^3)
		= \mathcal{O}(n^3)$.
	\end{proof}

	\begin{proposition}[Space Complexity of the Subgraph Algorithm]
		\label{prop:subgraph_space}
		The Subgraph Algorithm uses $\mathcal{O}(n^2)$ space for the two
		adjacency matrices and $\mathcal{O}(n)$ additional space for the
		signature arrays.
	\end{proposition}

	\begin{theorem}[Correctness of the Subgraph Algorithm]
		\label{thm:subgraph_correct}
		For any graphs $G$ and $G'$, the Subgraph Algorithm returns
		$\mathtt{true}$ if and only if $G \hookrightarrow G'$.
	\end{theorem}

	\begin{proof}
		Cyclic rotation of the column sequence of $B$ corresponds to a cyclic
		re-labelling of the vertices of $G'$.  Such a re-labelling preserves all
		edge adjacencies: if $(v_i, v_j) \in E_{G'}$, then the rotated columns
		still encode this edge under the new vertex names.  The $n_B$ rotations
		therefore cover $n_B$ distinct cyclic labellings of $G'$, forming a
		representative family for the structural similarity test.  By
		\Cref{lem:injective}, the injectivity of signatures ensures no spurious
		matches occur.  Hence the LCS test (\Cref{thm:subgraph_criterion})
		correctly certifies the presence or absence of a subgraph embedding.
	\end{proof}

	\subsection{Worked Example}

	\begin{example}[Subgraph Detection via Signatures]
		\label{ex:subgraph}
		Let $G$ be a directed path on four vertices with adjacency matrix
		\[
		A = \begin{pmatrix}
			0 & 1 & 0 & 0 \\
			0 & 0 & 1 & 0 \\
			0 & 0 & 0 & 1 \\
			0 & 0 & 0 & 0
		\end{pmatrix},
		\]
		and let $G'$ be the same path with one additional edge ($v_0 \to v_2$),
		\[
		B = \begin{pmatrix}
			0 & 1 & 1 & 0 \\
			0 & 0 & 1 & 0 \\
			0 & 0 & 0 & 1 \\
			0 & 0 & 0 & 0
		\end{pmatrix}.
		\]
		The Subgraph Algorithm computes:
		\begin{itemize}
			\item Signatures of $G$:
			$[\sigma_0^A, \sigma_1^A, \sigma_2^A, \sigma_3^A] = [1,\; 18,\; 36,\; 48]$;
			\quad row components: $\mathbf{r}^A = [1,\; 2,\; 4,\; 0]$.
			\item Signatures of $G'$:
			$[\sigma_0^B, \sigma_1^B, \sigma_2^B, \sigma_3^B] = [5,\; 18,\; 36,\; 48]$;
			\quad row components: $\mathbf{r}^B = [5,\; 2,\; 4,\; 0]$.
		\end{itemize}
		At rotation $k = 0$:
		$\mathrm{LCS}\bigl([1,\;2,\;4,\;0],\;[5,\;2,\;4,\;0]\bigr) = 3 \geq 2$,
		so the algorithm returns $\mathtt{true}$.  The containment
		$G \hookrightarrow G'$ is confirmed, and $G'$ is retained as the more
		informative graph.
	\end{example}

	%% ════════════════════════════════════════════════════════════════════════════
	\section{Algorithm Design}
	\label{sec:algo}
	
	\subsection{High-Level Overview}
	
	The HS-DFS algorithm proceeds in two phases:
	
	\begin{enumerate}
		\item \textbf{Meta-Graph Construction (MGC):} For every ordered pair
		$(G_i, G_j)$ with $i \neq j$, invoke the \textbf{Subgraph Algorithm}
		$\mathcal{S}$ (\Cref{alg:subgraph}) on the adjacency matrices $A_i$ and
		$A_j$.  If $\mathcal{S}(A_i, A_j) = \mathtt{true}$, insert the directed
		edge $(G_i \to G_j)$ into $\mathcal{M}$.
		\item \textbf{DFS Forest Computation (DFC):} Execute standard DFS on the
		constructed meta-graph $\mathcal{M}$, assigning discovery and finish
		timestamps and classifying all edges.
	\end{enumerate}
	
	\subsection{Pseudocode}
	
	\begin{algorithm}[H]
		\caption{HS-DFS: Hierarchical Subgraph Depth-First Search}
		\label{alg:hsdfs}
		\begin{algorithmic}[1]
			\Require Collection $\mathcal{G} = \{G_1, \ldots, G_N\}$ with adjacency
			matrices $\{A_1, \ldots, A_N\}$,
			Subgraph Algorithm oracle $\mathcal{S}$ (\Cref{alg:subgraph})
			\Ensure  Meta-graph $\mathcal{M}$, DFS forest $\mathcal{F}$,
			edge classification $\chi$
			
			\Statex \hrulefill
			\Statex \textbf{Phase 1 -- Meta-Graph Construction}
			\Statex \hrulefill
			
			\State $\mathcal{M} \leftarrow (\mathcal{G},\; \emptyset)$
			\LineComment{Initialise meta-graph with empty edge set}
			
			\For{$i \leftarrow 1$ \textbf{to} $N$}
			\For{$j \leftarrow 1$ \textbf{to} $N$}
			\If{$i \neq j$}
			\If{$\mathcal{S}(A_i,\, A_j) = \mathtt{true}$}
			\State $\mathcal{E} \leftarrow \mathcal{E} \cup \{(G_i, G_j)\}$
			\LineComment{Subgraph Algorithm (\Cref{alg:subgraph}): $G_i \hookrightarrow G_j$}
			\EndIf
			\EndIf
			\EndFor
			\EndFor
			\algstore{hsdfs}
		\end{algorithmic}
	\end{algorithm}

	\newpage

	\begin{algorithm}[H]
		\ContinuedFloat
		\caption{HS-DFS (continued)}
		\begin{algorithmic}[1]
			\algrestore{hsdfs}
			\Statex \hrulefill
			\Statex \textbf{Phase 2 -- DFS Forest Computation}
			\Statex \hrulefill
			
			\State $\mathit{colour}[G_i] \leftarrow \mathtt{WHITE}$ for all $i$
			\State $\mathit{parent}[G_i] \leftarrow \mathtt{NIL}$ for all $i$
			\State $\mathit{time} \leftarrow 0$
			
			\For{each $G_i \in \mathcal{G}$}
			\If{$\mathit{colour}[G_i] = \mathtt{WHITE}$}
			\State \Call{DFS-Visit}{$G_i$}
			\EndIf
			\EndFor
			
			\Statex
			\Procedure{DFS-Visit}{$u$}
			\State $\mathit{colour}[u] \leftarrow \mathtt{GRAY}$
			\State $\mathit{time} \leftarrow \mathit{time} + 1$
			\State $d[u] \leftarrow \mathit{time}$
			\LineComment{Discovery timestamp}
			\For{each $(u, v) \in \mathcal{E}$}
			\State $\chi(u,v) \leftarrow$ \Call{Classify-Edge}{$u$, $v$}
			\If{$\mathit{colour}[v] = \mathtt{WHITE}$}
			\State $\mathit{parent}[v] \leftarrow u$
			\State \Call{DFS-Visit}{$v$}
			\EndIf
			\EndFor
			\State $\mathit{colour}[u] \leftarrow \mathtt{BLACK}$
			\State $\mathit{time} \leftarrow \mathit{time} + 1$
			\State $f[u] \leftarrow \mathit{time}$
			\LineComment{Finish timestamp}
			\EndProcedure
			
			\Statex
			\Function{Classify-Edge}{$u$, $v$}
			\If{$\mathit{colour}[v] = \mathtt{WHITE}$}
			\Return $\mathtt{TREE}$
			\ElsIf{$\mathit{colour}[v] = \mathtt{GRAY}$}
			\Return $\mathtt{BACK}$
			\ElsIf{$d[u] < d[v]$}
			\Return $\mathtt{FORWARD}$
			\Else
			\Return $\mathtt{CROSS}$
			\EndIf
			\EndFunction
			
		\end{algorithmic}
	\end{algorithm}
	
	\subsection{Structural Justification for DFS over BFS}
	
	A critical design decision is the exclusive use of DFS in Phase~2.
	We formalise this below.
	
	\begin{theorem}[DFS Forest Property]
		\label{thm:dfsforest}
		DFS on any directed graph $\mathcal{M} = (\mathcal{G}, \mathcal{E})$
		produces a \emph{DFS forest} $\mathcal{F}$ partitioning $\mathcal{E}$
		into tree, back, forward, and cross edges.  BFS produces a \emph{BFS tree}
		(not a forest), yields no back or forward edge classification, and
		does not expose the full strongly-connected component structure.
	\end{theorem}
	
	\begin{proof}
		By the standard DFS classification theorem \cite[Ch.~22]{cormen2009}, every
		edge $(u,v)$ receives exactly one of the four labels based on the colour of
		$v$ when the edge is first explored and the relative order of discovery
		timestamps.  BFS assigns no finish times, cannot distinguish forward from
		cross edges, and therefore cannot determine whether $\mathcal{M}$ contains
		cycles -- which is essential for detecting circular containment hierarchies.
	\end{proof}
	
	\begin{corollary}
		The DFS forest of $\mathcal{M}$ directly encodes the \emph{subgraph
			containment hierarchy} of $\mathcal{G}$.  Tree edges represent direct
		containment; back edges indicate cyclic mutual containment; forward edges
		encode transitive containment shortcuts.
	\end{corollary}
	
	%% ════════════════════════════════════════════════════════════════════════════
	\section{Formal Complexity Analysis}
	\label{sec:complexity}
	
	Let $n$ denote a unified size parameter: each graph $G_i \in \mathcal{G}$ has
	at most $n$ vertices and $n^2$ edges, and the collection has $N = n$ members.
	
	\subsection{Phase 1: Meta-Graph Construction}
	
	\begin{theorem}[MGC Complexity]
		\label{thm:mgc}
		Phase~1 runs in time $\Theta(N^2 \cdot T_{\mathcal{S}}(n))$.
		With the Subgraph Algorithm as oracle
		($T_{\mathcal{S}}(n) = \mathcal{O}(n^3)$, \Cref{thm:subgraph_complexity})
		and $N = n$, this yields $\mathcal{O}(n^5)$.
	\end{theorem}
	
	\begin{proof}
		The double loop over $(i,j)$ with $i \neq j$ executes exactly
		$N(N-1)$ invocations of the Subgraph Algorithm.
		By \Cref{thm:subgraph_complexity}, each invocation costs
		$\mathcal{O}(n^3)$.  Substituting $N = n$:
		\[
		T_{\text{MGC}} = N(N-1) \cdot T_{\mathcal{S}}(n)
		= n(n-1) \cdot \mathcal{O}(n^3)
		= \mathcal{O}(n^5).
		\]
		The lower bound $\Omega(n^5)$ follows because for a dense meta-graph
		nearly all $N^2$ oracle calls are required.
	\end{proof}
	
	\subsection{Phase 2: DFS Forest Computation}
	
	\begin{theorem}[DFC Complexity]
		\label{thm:dfc}
		Phase~2 runs in time $\Theta(N + |\mathcal{E}|)$.  In the worst case
		$|\mathcal{E}| = N(N-1) = \mathcal{O}(N^2) = \mathcal{O}(n^2)$,
		giving $\Theta(n^2)$.
	\end{theorem}
	
	\begin{proof}
		Standard DFS processes each vertex once (cost $\mathcal{O}(N)$) and
		each edge once (cost $\mathcal{O}(|\mathcal{E}|)$).  With $N = n$ and
		at most $N^2$ meta-edges, $T_{\text{DFC}} = \mathcal{O}(n + n^2)
		= \mathcal{O}(n^2)$.
	\end{proof}
	
	\subsection{Total Complexity}
	
	\begin{theorem}[HS-DFS Total Complexity]
		\label{thm:total}
		The HS-DFS algorithm runs in time $\mathcal{O}(n^5)$ for dense
		meta-graphs, and in time $\mathcal{O}(n^4)$ for sparse meta-graphs
		where $|\mathcal{E}| = \mathcal{O}(n)$.
	\end{theorem}
	
	\begin{proof}
		\[
		T_{\text{HS-DFS}} = T_{\text{MGC}} + T_{\text{DFC}}
		= \mathcal{O}(n^5) + \mathcal{O}(n^2)
		= \mathcal{O}(n^5).
		\]
		In the sparse case, Phase~1 requires $\mathcal{O}(n) \cdot \mathcal{O}(n^3)
		= \mathcal{O}(n^4)$ oracle calls and Phase~2 costs $\mathcal{O}(n)$.
	\end{proof}
	
	\begin{table}[H]
		\centering
		\caption{Complexity summary by meta-graph density.}
		\label{tab:complexity}
		\begin{tabular}{lccc}
			\toprule
			\textbf{Scenario} & \textbf{Meta-edges} & \textbf{Phase~1} & \textbf{Phase~2} \\
			\midrule
			Dense (worst case)   & $\mathcal{O}(n^2)$ & $\mathcal{O}(n^5)$ & $\mathcal{O}(n^2)$ \\
			Average case         & $\mathcal{O}(n^{1.5})$ & $\mathcal{O}(n^{4.5})$ & $\mathcal{O}(n^{1.5})$ \\
			Sparse               & $\mathcal{O}(n)$   & $\mathcal{O}(n^4)$ & $\mathcal{O}(n)$ \\
			\textbf{Total (dense)} & -- & \multicolumn{2}{c}{$\boldsymbol{\mathcal{O}(n^5)}$} \\
			\bottomrule
		\end{tabular}
	\end{table}
	
	\subsection{Space Complexity}
	
	\begin{proposition}[Space Complexity]
		HS-DFS requires $\mathcal{O}(n^2 + n \cdot n^2) = \mathcal{O}(n^3)$ space:
		$\mathcal{O}(n^2)$ for the meta-graph adjacency matrix and
		$\mathcal{O}(n^3)$ for storing the $n$ input graphs of size $\mathcal{O}(n^2)$
		each.
	\end{proposition}
	
	\subsection{Lower Bound}
	
	\begin{theorem}[Conditional Lower Bound]
		Under the Strong Exponential Time Hypothesis (SETH), no algorithm can
		solve Phase~1 in time $o(n^5)$ in the general case.
	\end{theorem}
	
	\begin{proof}[Proof sketch]
		Subgraph isomorphism for general graphs requires $\Omega(n^3)$ per pair
		under SETH-based reductions \cite{williams2018}.  With $\Omega(n^2)$
		pairs in the dense case, the total lower bound is $\Omega(n^5)$.
	\end{proof}
	
	\subsection{Optimisation Opportunities}
	
	While the worst-case bound is tight, several practical optimisations reduce
	constant factors significantly:
	
	\begin{enumerate}[label=(\roman*)]
		\item \textbf{Size filtering:} Pairs $(G_i, G_j)$ with $n_i > n_j$
		cannot satisfy $G_i \hookrightarrow G_j$ and are rejected by the first
		check of the Subgraph Algorithm in $\mathcal{O}(1)$.
		\item \textbf{Signature pre-filtering:} The column-signature arrays of
		all graphs can be pre-computed once in $\mathcal{O}(N \cdot n^2)$.
		A quick set-intersection test on the row-component arrays
		($\mathcal{O}(n \log n)$) provides a necessary condition for subgraph
		containment, eliminating many oracle calls at negligible cost.
		\item \textbf{Rotation early termination:} The cyclic rotation loop in
		Step~3 of the Subgraph Algorithm terminates as soon as
		$\mathrm{LCS} \geq 2$ is found, often in far fewer than $n_B$ iterations.
		\item \textbf{Parallelism:} All $N^2$ Subgraph Algorithm invocations are
		mutually independent and can be distributed across $P$ processors,
		reducing Phase~1 to $\mathcal{O}(n^5 / P)$.
	\end{enumerate}
	
	%% ════════════════════════════════════════════════════════════════════════════
	\section{DFS Forest Structure and its Semantics}
	\label{sec:forest}
	
	The DFS forest $\mathcal{F}$ produced in Phase~2 carries rich semantic
	information that goes far beyond mere path-finding:
	
	\begin{description}
		\item[Tree edges $(u \to v)$:] $G_u$ is directly subgraph-isomorphic to
		$G_v$, and $G_v$ was first discovered through $G_u$.  These form the
		\emph{primary containment hierarchy}.
		
		\item[Back edges $(u \to v)$:] $G_v$ is an ancestor of $G_u$ in the DFS
		tree, implying a cycle in the containment relation.  This detects
		\emph{mutual or circular containment}, a structurally significant
		phenomenon.
		
		\item[Forward edges $(u \to v)$:] $G_v$ is a descendant of $G_u$ but not a
		direct child.  This encodes \emph{transitive containment}.
		
		\item[Cross edges $(u \to v)$:] $G_v$ belongs to a different DFS subtree.
		These reveal \emph{lateral containment} across independently evolved
		substructures.
		
		\item[Finish time ordering:] The reverse of finish-time ordering ($f[u]$)
		yields a \emph{topological sort} of the meta-graph if it is acyclic,
		representing a valid total order of containment depth.
	\end{description}
	
	%% ════════════════════════════════════════════════════════════════════════════
	\section{Application Domains}
	\label{sec:applications}
	
	The HS-DFS algorithm is not domain-specific: its inputs are graphs and its
	output is a structured hierarchy.  This generality makes it a powerful
	unifying primitive.  We now provide an extensive analysis of its applicability
	across nine major domains.
	
	%% ──────────────────────────────────────────────────────────────────────────
	\subsection{Neuroscience and Brain Research}
	\label{sec:brain}
	
	The human brain is among the most complex graph-structured systems known.
	At every spatial scale -- from synaptic connections between individual neurons,
	to local circuits, cortical columns, functional areas, and whole-brain
	networks -- the brain is hierarchically composed of sub-networks.
	The dedicated brain-research monograph \cite{epp2026brain} provides
	rich biological grounding for the HS-DFS application in neuroscience and
	motivates the extensions developed in the paragraphs below.
	
	\paragraph{The Connectome as a Graph of Graphs.}
	The \emph{connectome} is the complete map of neural connections in a brain.
	Modern neuroimaging (fMRI, DTI tractography, electron microscopy) produces
	multi-scale connectome data.  Each functional region (e.g., the default mode
	network, the visual processing hierarchy, the hippocampal memory circuit) can
	be represented as a graph $G_i$.  HS-DFS applied to the collection of all
	such regional graphs $\mathcal{G} = \{G_{\text{DMN}}, G_{\text{visual}},
	G_{\text{hippo}}, \ldots\}$ constructs the meta-graph $\mathcal{M}$ encoding
	which circuits are embedded within which larger circuits \cite{sporns2011}.
	
	\paragraph{The Invisible Life Force: Synaptic Transmission as Graph
		Traversal.}
	The fundamental carrier of information in the brain's hierarchical graph is
	the electro-chemical impulse flow from synapse to synapse -- what
	\cite{epp2026brain} calls the \emph{invisible life force} (unsichtbarer Geist
	des Lebens).  Every traversal of a meta-graph edge $(G_i \to G_j)$ in
	$\mathcal{M}$ corresponds, at the biological level, to an avalanche of action
	potentials propagating through the sub-network $G_i$ and activating circuits
	in $G_j$.  The elementary impulse is described by the Hodgkin--Huxley
	equations \cite{hodgkin1952}:
	\[
	C_m\,\frac{dV}{dt} = I_{\mathrm{Na}} + I_{\mathrm{K}} + I_L + I_{\mathrm{ext}},
	\]
	and synaptic transmission is quantised \cite{delcastillo1954}: the mean
	postsynaptic potential amplitude is
	\[
	\mu_{\mathrm{EPSP}} = n \cdot p \cdot q,
	\]
	where $n$ is the number of release sites, $p$ the vesicle-release
	probability, and $q$ the quantal amplitude.  The \emph{life-force intensity
		index} \cite{epp2026brain}
	\[
	\mathcal{G}(\mathcal{N},\,t) \;:=\;
	\sum_{s \in \mathcal{S}} \lambda_s(t)\cdot w_s \cdot \phi_s
	\]
	(action-potential rate $\lambda_s$, synaptic weight $w_s$, transmission
	efficiency $\phi_s$) quantifies the aggregate synaptic transmission rate and
	serves as the biological \emph{edge-weight function} of $\mathcal{M}$: edges
	$(G_i \to G_j)$ with low $\mathcal{G}$ are weakly activated links in the
	containment hierarchy.  Neurotransmitters are the specific carriers of
	distinct life-force aspects: glutamate drives long-term potentiation (LTP)
	and NMDA-dependent truth search; dopamine boosts prefrontal signal-to-noise
	ratio; acetylcholine gates thalamo-cortical arousal; GABAergic
	inhibition under the epistemic cloud suppresses all three simultaneously
	\cite{epp2026brain}.
	
	\paragraph{The Epistemic Cloud: Neurobiological Obstruction of Graph Access.}
	\cite{epp2026brain} formally defines the \emph{epistemic cloud}
	$\mathcal{W}$ with cloud density $\delta(\mathcal{W}) \in [0,1]$:
	\[
	\delta(\mathcal{W}) \;:=\; 1 - \frac{|\mathcal{I}_P \cap \mathcal{I}^*|}{|\mathcal{I}^*|},
	\]
	where $\mathcal{I}_P$ is the person's available information and
	$\mathcal{I}^*$ the complete information required for optimal problem
	solving.  Biologically, the epistemic cloud manifests as altered prefrontal
	connectivity, suppressed gamma coherence, elevated amygdala reactivity, and
	chronically elevated cortisol levels -- all neurobiologically measurable
	with EEG and fMRI \cite{epp2026brain}.  Cortisol reduces synaptic density
	via dendritic retraction in PFC pyramidal cells \cite{radley2004} and
	suppresses the life force:
	\[
	\frac{d\mathcal{G}}{d\delta} \;\leq\; 0.
	\]
	In HS-DFS terms, the epistemic cloud corresponds to a \emph{degraded
		meta-graph}: edges are removed or weakened, the DFS forest becomes
	shallower and sparser, and correct hierarchical containment relationships
	can no longer be detected.  The cloud obstructs subgraph containment
	detection at the oracle level: dopaminergic hypofuntion in the dlPFC
	\cite{goldmanrakic1995} or over-dominant GABAergic inhibition may prevent
	the Subgraph Algorithm from detecting the embedding
	$G_i \hookrightarrow G_j$.  Accordingly, \cite{epp2026brain} proves
	a \emph{double-blockade theorem}: the epistemic cloud degrades the
	truth-access index super-linearly because it simultaneously raises
	$\delta$ directly \emph{and} suppresses $\mathcal{G}$,
	\[
	\Delta\mathrm{WZI} \;\geq\; \delta + \varepsilon(\delta),\quad \varepsilon > 0,
	\]
	with the \emph{truth-access index} \cite{epp2026brain}
	\[
	\mathrm{WZI}(P, \mathcal{G}, \mathcal{W}) \;:=\;
	\bigl(1 - \delta(\mathcal{W})\bigr) \cdot
	\frac{\mathcal{G}(\mathcal{N},t)}{\mathcal{G}_{\max}} \cdot
	\mathbf{1}_{\omega(t) < 0},
	\]
	where $\mathbf{1}_{\omega < 0}$ is the clockwise-rotation indicator (see
	below).  WZI\,$= 1$ corresponds to a fully traversable meta-graph;
	WZI\,$= 0$ to complete blockade.
	
	\paragraph{Clockwise Cortical Information Processing (Rechtsdrehung).}
	A central finding of \cite{epp2026brain} is that cortical information
	processing proceeds in a \emph{clockwise direction} when viewed from
	dorsal: the activation wave propagates anterior $\to$ lateral-right $\to$
	posterior $\to$ lateral-left $\to$ anterior, with instantaneous angular
	velocity $\omega(t) < 0$ (Definition~3.3 in \cite{epp2026brain}):
	\[
	\omega(t) \;=\;
	\frac{A_x(t)\,\dot{A}_y(t) - A_y(t)\,\dot{A}_x(t)}{|\vec{A}(t)|^2}
	\;<\; 0,
	\]
	where $\vec{A}(t) = (A_x(t), A_y(t))$ is the cortical activation vector.
	Six independent criteria confirm this clockwise rotation in
	\cite{epp2026brain}: (i)~negative instantaneous frequency
	$\langle\omega\rangle < 0$ from the Hilbert-transform analytic signal of
	the Gamma band; (ii)~Phase-Amplitude Coupling (PAC) peak at theta phase
	$-\pi/2$, shifting to $+\pi/3$ under the epistemic cloud; (iii)~the
	phase-lead matrix: frontal electrodes consistently lead posterior ones;
	(iv)~the polar hodograph winding clockwise in phase space; (v)~the
	rotation tensor showing $\omega < 0$ stably at $\approx{-40}$\,Hz;
	(vi)~Fourier analysis of the gamma peak confirming negative phase shift.
	Under the epistemic cloud, the Default Mode Network (DMN) gains
	dominance, gamma coherence collapses, and the clockwise rotation
	destabilises or inverts \cite{epp2026brain}.
	
	The clockwise rotation has a direct graph-theoretic interpretation for
	HS-DFS: the DFS traversal order of the meta-graph mirrors the cortical
	activation cycle.  A complete clockwise cycle of $\approx 25\,\mathrm{ms}$
	at $|\omega_0| \approx 40\,\mathrm{Hz}$ corresponds to one complete DFS
	pass through all directly reachable sub-networks.  The term
	$\mathbf{1}_{\omega < 0}$ in the WZI formula acts as a binary gate:
	a brain not operating in the clockwise mode cannot complete the
	meta-graph traversal reliably.  Theorem~3.6 of \cite{epp2026brain}
	(Phase-Lead Theorem) states that in the clockwise cortex the frontal
	electrode leads the posterior one ($\Delta\varphi_{ij} > 0$) and
	this sign reverses under the epistemic cloud -- a directly testable
	prediction.
	
	\paragraph{Ad-hoc Information Selection via the Thalamo-Cortical Gateway.}
	Memory in the brain is not a passive store but an \emph{on-demand
		retrieval system}: the brain fetches from its vast repository only
	those contents relevant at the moment of query, steered by prefrontal
	executive control and the ARAS thalamo-cortical gateway
	\cite{epp2026brain}.  This \emph{ad-hoc information provision} (AIB) is
	modelled as
	\[
	\mathcal{A}(q, \mathcal{M})
	\;:=\; \arg\max_{I \subseteq \mathcal{M}}\;
	\mathrm{Rel}(I,\,q) \cdot \mathrm{Avail}(I,\,\mathcal{G}),
	\]
	where $\mathrm{Avail}(I,\mathcal{G}) = \sigma((\mathcal{G} -
	\mathcal{G}_\theta)/\tau_\mathcal{G})$ is a sigmoid gate that opens
	fully only when $\mathcal{G}$ exceeds threshold $\mathcal{G}_\theta$
	(Theorem~3.7 in \cite{epp2026brain}).  The selectivity efficiency is
	bounded by
	\[
	\eta_{\mathcal{A}} \;\leq\;
	\frac{\mathcal{G}}{\mathcal{G}_{\max}}\cdot(1-\delta) \;=\;\mathrm{WZI}.
	\]
	Reaction-time simulations in \cite{epp2026brain} (Experiment~8) show
	that the ad-hoc retrieval latency rises from $\approx 320$\,ms (clear
	life force, $\delta \approx 0$) to $\approx 840$\,ms (dense cloud,
	$\delta \approx 0.8$), with the ARAS gateway flickering chaotically at
	intermediate cloud densities.
	
	In HS-DFS terms, each call to \textsc{DFS-Visit}$(u)$ models the
	thalamo-cortical query that retrieves which sub-networks $v$ are
	reachable from $u$.  Under the epistemic cloud the gateway flickers and
	some edges $(u \to v)$ are missed; the resulting DFS forest is sparser
	than $\mathcal{M}^*$, and containment hierarchies go undetected.
	
	\paragraph{Cognitive Wetting: A Prerequisite for Neural Hierarchy
		Traversal.}
	\cite{epp2026brain} identifies a fifth dimension of cognition --
	\emph{cognitive wetting} ($\mathcal{B}$) -- defined as the fraction of
	synaptic connections uniformly permeated by active electro-chemical
	information flow:
	\[
	\mathcal{B}(t) \;:=\; \frac{1}{N}\sum_{i=1}^{N}
	\mathbf{1}\bigl[\Phi_i(t) \geq \Phi_{\min}\bigr] \;\in\; [0,1].
	\]
	The epistemic cloud reduces wetting super-linearly \cite{epp2026brain}
	via three concurrent mechanisms: (i)~\emph{transmitter depletion}
	through cortisol-mediated suppression of dopamine and serotonin
	biosynthesis; (ii)~\emph{membrane hyperpolarisation} through elevated
	GABAergic inhibition; (iii)~\emph{glial dysfunction} through astrocyte
	shrinkage and resulting $\mathrm{K}^+$ homeostasis failure.  The formal
	result \cite{epp2026brain}:
	\[
	\mathcal{B}(\delta) \;=\; \mathcal{B}_0 \cdot (1-\delta)^\alpha,
	\quad \alpha \approx 1.8,
	\]
	yields a critical percolation threshold $\mathcal{B}^* \approx 0.75$
	below which cognitive capacity collapses non-linearly -- a
	phase-transition-like breakdown established via a Hopfield network
	capacity argument \cite{hopfield1982,epp2026brain}:
	\[
	\frac{dC_{\max}}{d\mathcal{B}}\bigg|_{\mathcal{B}=\mathcal{B}^*}
	\;=\; -\infty \quad \text{(thermodynamic limit)}.
	\]
	In the HS-DFS context, cognitive wetting determines the \emph{graph
		connectivity of $\mathcal{M}$}: below $\mathcal{B}^*$ the meta-graph
	undergoes a percolation transition and fragments into disconnected
	components, the DFS forest decomposes into many isolated trees, and
	the hierarchical containment structure collapses.
	
	\paragraph{The Integrated Formal Model of Cognition.}
	\cite{epp2026brain} unifies all five dimensions into a single equation
	for free, unconstrained cognition:
	\[
	\Psi_{\mathrm{free}}(t) \;=\;
	\underbrace{\mathcal{B}(\delta, T, V_m)}_{\text{wetting}}
	\;\cdot\;
	\underbrace{\mathrm{IZI}(t)}_{\text{life force}}
	\;\cdot\;
	\underbrace{(1-\delta)}_{\text{cloud freedom}}
	\;\cdot\;
	\underbrace{\mathbf{1}_{\omega(t)<0}}_{\text{clockwise mode}},
	\]
	with threshold $\Psi_{\mathrm{free}} \geq \Psi^* \approx 0.6$.  All
	four factors are multiplicative: if any one collapses to zero, free
	cognition is impossible regardless of the others.  The epistemic cloud
	directly degrades \emph{two} factors simultaneously ($\mathcal{B}$ via
	Theorem on Wetting Obstruction and $(1-\delta)$ directly), making it the
	most potent single inhibitor of hierarchical graph processing in the brain.

	For HS-DFS, the integrated model translates into a \emph{biological
		feasibility condition}: the brain can reliably apply a subgraph oracle
	(i.e., HS-DFS operates correctly at the neural level) if and only if
	$\Psi_{\mathrm{free}} \geq \Psi^*$.  Bayesian update experiments in
	\cite{epp2026brain} (Experiment~3) show that a cloud-affected mind
	plateaus at only $\approx 62.5\,\%$ of the truth asymptotically, even
	given unlimited evidence -- a direct analogue of an HS-DFS oracle that
	permanently misses $37.5\,\%$ of subgraph containment edges.
	
	\paragraph{Clinical Applications.}
	The DFS forest of the brain's meta-graph provides a \emph{hierarchical
		fingerprint} of neural organisation.  Deviations from the healthy forest
	structure -- detected by comparing patient and reference forests -- provide
	diagnostic markers for neurological disorders.  The biological framework of
	\cite{epp2026brain} extends and sharpens these markers:
	\begin{itemize}
		\item \textbf{Alzheimer's disease:} Progressive degradation of the
		hippocampal sub-network manifests as the loss of edges in $\mathcal{M}$
		incident to $G_{\text{hippo}}$, fragmenting the DFS tree.  In the
		language of \cite{epp2026brain}, Alzheimer's is a \emph{progressive
			cognitive dehydration}: a $30$--$60\,\%$ synapse loss drives
		$\mathcal{B} \to 0$, causing meta-graph fragmentation before
		memory content is lost.
		\item \textbf{Epilepsy:} Abnormal back edges in the DFS forest indicate
		pathological feedback loops (circular containment) in the seizure
		focus network, corresponding biologically to uncontrolled re-entrant
		activation cycles in the clockwise cortical wave.
		\item \textbf{Schizophrenia:} Anomalous cross edges reveal atypical
		lateral connections between normally segregated networks.
		\item \textbf{Depression:} Following \cite{epp2026brain}, major
		depression features elevated cortisol, reduced monoaminergic
		transmission, and increased GABAergic activity -- precisely the
		parameters that maximally reduce $\mathcal{B}$ and suppress
		$\mathcal{G}$.  The DFS forest of a depressed connectome is
		structurally shallower, consistent with cognitive rigidity and
		reduced associative thinking; EEG instantaneous frequency
		$\langle\omega\rangle$ approaches zero (loss of clockwise rotation),
		observable without invasive measurements.
		\item \textbf{EEG biomarker:} \cite{epp2026brain} proposes a
		standardised 20-minute EEG protocol yielding an \emph{epistemic
			cloud score}: $\langle\omega\rangle < 0$ = healthy clockwise
		processing; $\langle\omega\rangle \to 0$ or $> 0$ = cloud indicator.
		The PAC peak angle (healthy: $-\pi/2$; cloud: $+\pi/3$) and the
		phase-lead matrix structure (six-criterion scorecard in
		\cite{epp2026brain}) provide independent confirmation and could
		serve as objective biomarkers for therapeutic success.
	\end{itemize}
	
	\paragraph{Consciousness Research.}
	Theories of consciousness (Integrated Information Theory, Global Workspace
	Theory) posit that conscious states correspond to globally integrated
	information patterns.  The DFS forest directly quantifies integration: the
	depth and breadth of the tree reflect the degree to which sub-circuits are
	embedded within broader structures, providing a computable proxy for
	$\Phi$ (phi), the IIT measure of integrated information \cite{iit}.
	\cite{epp2026brain} adds a precise biological complement: cognitive wetting
	$\mathcal{B}$ is the necessary physical correlate of the \emph{unity of
		consciousness} -- as $\mathcal{B} \to 0$ the cortical network fragments
	into cognitive islands without mutual communication, phenomenologically
	corresponding to states of deep sleep, deep anaesthesia, or certain
	dissociative disorders.
	
	\paragraph{Neural Development and Transgenerational Effects.}
	During brain maturation, sub-circuits form sequentially.  HS-DFS applied
	longitudinally (at multiple developmental time points) tracks how the
	meta-graph evolves, detecting when new subgraph relationships emerge,
	corresponding to developmental milestones such as language acquisition
	or the maturation of executive function.  Crucially, \cite{epp2026brain}
	demonstrates via an epidemiological SI-model ($R_0 = \beta/\gamma = 4$)
	and transgenerational epigenetics analysis that the epistemic cloud
	propagates across generations: perinatal glucocorticoids programme the
	foetal HPA axis and alter FKBP5/NR3C1 methylation patterns in the
	child, pre-programming offspring for a cloud state from birth.  In the
	HS-DFS framework this means that the meta-graph of an affected brain
	is structurally impoverished from the outset -- not through a lack of
	neural connections, but through systematic suppression of the subgraph
	oracle's sensitivity.  Educational interventions reduce cloud density
	across five simulated generations (Experiment~4 in \cite{epp2026brain}),
	corresponding to a gradual restoration of the meta-graph's edge density
	and DFS tree depth.
	
	%% ──────────────────────────────────────────────────────────────────────────
	\subsection{Genomics and Systems Biology}
	\label{sec:genomics}
	
	\paragraph{Gene Regulatory Networks.}
	A gene regulatory network (GRN) $G = (V, E)$ has genes as vertices and
	directed regulatory interactions (activations, repressions) as edges.
	Organisms contain hundreds of functional GRN sub-modules: the cell cycle
	module, the apoptosis module, the p53 pathway, the Wnt signalling cascade,
	and so forth.
	
	Applying HS-DFS to the collection of all known pathway graphs in a model
	organism (e.g., \textit{Homo sapiens} KEGG database) constructs the
	\emph{pathway meta-graph}.  Subgraph relationships here correspond to
	\emph{pathway nesting}: smaller signalling motifs (negative feedback loops,
	feed-forward loops) are structurally embedded within larger pathways.
	
	\paragraph{Motif Hierarchy Discovery.}
	Network motifs -- recurring subgraph patterns in biological networks -- form
	a natural hierarchy.  A 3-node feedback loop is subgraph-isomorphic to a
	5-node regulatory module which is in turn embedded in a full pathway.  The
	DFS forest produced by HS-DFS is precisely the \emph{motif hierarchy}, a
	previously difficult structure to compute efficiently.
	
	\paragraph{Cancer Genomics.}
	Tumour cells exhibit dysregulated GRNs.  Comparing the meta-graph of a
	healthy cell's GRN collection to that of a tumour cell exposes which pathway
	sub-modules have been corrupted (missing edges in $\mathcal{M}$) or
	aberrantly fused (new spurious edges).  The DFS forest difference directly
	points to candidate therapeutic targets.
	
	\paragraph{Comparative Genomics.}
	Across species, orthologous pathways are graph-isomorphic variants of the
	same biological process.  HS-DFS on the union of human and mouse pathway
	graphs quantifies evolutionary conservation at the subgraph level, revealing
	which regulatory modules are deeply conserved (appear as high-confidence
	tree edges in the meta-graph) versus species-specific (isolated nodes or
	cross edges between evolutionary lineages).
	
	\paragraph{Protein Interaction Networks.}
	Protein complexes and interaction sub-networks exhibit similar hierarchical
	structure.  Small protein interaction motifs (hub-and-spoke, cliques) are
	subgraph-isomorphic to larger complexes, which are in turn contained in
	cellular machinery modules.  HS-DFS on the proteome-wide interaction graph
	collection provides a \emph{proteome containment map}.
	
	%% ──────────────────────────────────────────────────────────────────────────
	\subsection{Logistics and Supply Chain Management}
	\label{sec:logistics}
	
	\paragraph{Multi-Level Supply Networks.}
	Modern supply chains are hierarchically structured: individual factory
	production graphs are embedded within regional distribution graphs, which
	are in turn embedded within continental logistics networks.  Formally, each
	level is a graph $G_i$ and the collection $\mathcal{G}$ is the full
	multi-tier supply network.
	
	HS-DFS constructs the meta-graph encoding which regional sub-networks are
	structurally contained within global network models, enabling:
	
	\begin{itemize}
		\item \textbf{Bottleneck detection:} Nodes $G_i$ with high in-degree in
		$\mathcal{M}$ are \emph{critical hubs}: many other sub-networks embed
		within them, meaning their disruption cascades globally.
		\item \textbf{Redundancy analysis:} Multiple parallel subgraph embeddings
		of the same $G_i$ indicate that the logistics function is replicated,
		providing resilience.
		\item \textbf{Reorganisation planning:} The topological sort of $\mathcal{M}$
		(from the DFS finish times) gives a dependency-respecting order for
		rolling out infrastructure changes.
	\end{itemize}
	
	\paragraph{Route Optimisation.}
	In last-mile delivery, city districts are modelled as local route graphs.
	A vehicle routing plan for a district is subgraph-isomorphic to the full
	city graph.  HS-DFS identifies which district plans can be executed
	simultaneously without interference (cross edges in the meta-graph) versus
	which must be sequenced (tree edges encoding dependencies).
	
	\paragraph{Inventory and Warehouse Networks.}
	Warehouse flow graphs (items flowing between storage zones) can be compared
	across facilities.  HS-DFS on a collection of warehouse graphs identifies
	facilities that share structural flow patterns, enabling the transfer of
	optimisation policies between facilities with matching subgraph structure.
	
	%% ──────────────────────────────────────────────────────────────────────────
	\subsection{Industrial Automation and Cyber-Physical Systems}
	\label{sec:automation}
	
	\paragraph{Petri Net Composition.}
	Industrial processes are commonly modelled as Petri nets or timed automata.
	A manufacturing cell's process model is a graph; an assembly line consists
	of multiple such cells; a factory floor is a meta-graph of assembly lines.
	HS-DFS on the collection of all process models in a factory detects which
	sub-process models are structurally embedded within larger models, enabling
	\emph{modular process reuse}: if $G_{\text{cell-A}} \hookrightarrow
	G_{\text{line-3}}$, the control logic of cell~A can be safely re-deployed
	as part of line~3 without redesign.
	
	\paragraph{AUTOSAR and Automotive Software Architecture.}
	In AUTOSAR-based automotive systems, software components (SWCs) are connected
	through ports and communication links, forming a directed graph.  Each
	vehicle function (engine management, braking, ADAS) is a sub-graph.
	HS-DFS on the AUTOSAR component collection enables:
	\begin{itemize}
		\item Detecting which SWC configurations from one ECU are
		subgraph-isomorphic to configurations on another ECU
		(enabling \emph{cross-ECU reuse}).
		\item Finding all functional sub-systems that must be preserved during
		an OTA software update by identifying their containment relationships.
	\end{itemize}
	
	\paragraph{Robotics and Task Graph Planning.}
	Robot task plans are directed acyclic graphs (DAGs) over atomic actions.
	Complex manipulation tasks contain sub-tasks that re-appear across different
	high-level plans.  HS-DFS on the library of robot task graphs constructs the
	\emph{task containment hierarchy}, enabling a robot controller to reuse
	solutions to sub-tasks instead of re-planning from scratch -- a form of
	\emph{macro-action learning} grounded in formal subgraph structure.
	
	\paragraph{Industrial IoT Topology Management.}
	In large-scale IoT deployments (smart factories, smart grids), device
	communication graphs are hierarchically composed.  Sensor clusters form
	local graphs; edge computing zones aggregate clusters; cloud management
	spans all zones.  HS-DFS provides a principled mechanism for:
	\begin{itemize}
		\item Firmware update sequencing (DFS finish-time order).
		\item Fault propagation analysis (back edges indicate feedback loops
		in control messages that may amplify failures).
		\item Network slicing validation (checking that a proposed network
		slice graph is subgraph-isomorphic to the physical network graph).
	\end{itemize}
	
	%% ──────────────────────────────────────────────────────────────────────────
	\subsection{Software Engineering and Program Analysis}
	\label{sec:software}
	
	\paragraph{Call Graph and Dependency Analysis.}
	A software system's call graph $G_{\text{sys}}$ contains the call graphs of
	all its modules as induced sub-graphs.  HS-DFS on the collection of all
	module call graphs in a large code base detects which modules are structurally
	contained within which subsystems.  This enables:
	\begin{itemize}
		\item \textbf{Dead code elimination:} Modules not appearing as subgraphs in
		any higher-level graph are candidates for removal.
		\item \textbf{Refactoring guidance:} Modules appearing as subgraphs in
		multiple parent modules are candidates for extraction into shared
		libraries.
		\item \textbf{Dependency cycle detection:} Back edges in the DFS forest
		reveal circular module dependencies, a known code quality defect.
	\end{itemize}
	
	\paragraph{Microservice Architecture.}
	Microservice communication graphs evolve over time.  HS-DFS on versioned
	snapshots of the service graph collection detects structural regressions:
	a new deployment has broken subgraph containment that existed in the previous
	version.
	
	\paragraph{Type Hierarchy in Object-Oriented Systems.}
	Class inheritance hierarchies are graphs.  HS-DFS on the collection of class
	interaction graphs in an OO system validates Liskov Substitution Principle
	compliance: if $G_{\text{child}} \hookrightarrow G_{\text{parent}}$
	(the child's interaction graph is subgraph-isomorphic to the parent's),
	the child is a valid structural substitute.
	
	%% ──────────────────────────────────────────────────────────────────────────
	\subsection{Cheminformatics and Drug Discovery}
	\label{sec:chemo}
	
	\paragraph{Molecular Subgraph Isomorphism.}
	Molecules are naturally represented as graphs (atoms as vertices, bonds as
	edges).  Pharmacophores -- the functionally active sub-structures of drug
	molecules -- are subgraphs.  HS-DFS on a library of molecular graphs
	$\mathcal{G} = \{G_{\text{mol}_1}, \ldots, G_{\text{mol}_N}\}$ produces:
	\begin{itemize}
		\item The \emph{pharmacophore hierarchy}: which pharmacophores are embedded
		within which drug candidates.
		\item A \emph{scaffold tree}: the DFS forest of the molecular library,
		organising compounds by shared substructure, directly useful for
		scaffold-hopping in medicinal chemistry.
		\item \emph{ADMET property propagation}: if compound $G_i$ is a subgraph
		of $G_j$, and $G_i$ has known toxicity, HS-DFS propagates this
		risk annotation upward through the DFS tree.
	\end{itemize}
	
	\paragraph{Reaction Network Analysis.}
	Chemical reaction networks (CRNs) model biochemical or industrial reactions
	as hypergraphs.  HS-DFS on collections of CRN sub-networks identifies shared
	reaction motifs, supporting the design of novel catalytic processes by
	reusing known sub-network solutions.
	
	%% ──────────────────────────────────────────────────────────────────────────
	\subsection{Knowledge Graphs and Semantic Web}
	\label{sec:kg}
	
	\paragraph{Ontology Alignment.}
	Knowledge graphs such as Wikidata, DBpedia, or domain ontologies (Gene
	Ontology, Medical Subject Headings) are graph-structured.  HS-DFS on a
	collection of domain-specific sub-ontologies detects which ontologies are
	structurally compatible (subgraph-isomorphic), enabling automatic ontology
	alignment and merging -- a long-standing challenge in the Semantic Web.
	
	\paragraph{Query Optimisation.}
	SPARQL queries over knowledge graphs are themselves graph patterns.  A
	SPARQL query $Q$ can be evaluated as a subgraph isomorphism problem:
	find all embeddings of $Q$ in the KG.  HS-DFS on a collection of frequently
	used SPARQL queries identifies which queries are subgraph-isomorphic to
	others, enabling query result caching at the structural level.
	
	%% ──────────────────────────────────────────────────────────────────────────
	\subsection{Social Network Analysis}
	\label{sec:social}
	
	\paragraph{Community Structure and Overlapping Groups.}
	Social networks decompose into communities, which are sub-graphs.  HS-DFS
	on the collection of detected communities reveals \emph{community nesting}:
	small groups that are structurally embedded within larger ones.  This is
	relevant for:
	\begin{itemize}
		\item Influence propagation modelling (tree edges in the meta-graph
		are primary influence channels).
		\item Disinformation spread analysis (back edges indicate echo-chamber
		feedback loops).
		\item Recommendation systems (users in communities that are cross-edge-linked
		share interests via lateral subgraph similarity).
	\end{itemize}
	
	\paragraph{Organisational Network Analysis.}
	In enterprises, formal reporting structures and informal communication
	networks are graphs.  HS-DFS on departmental communication graphs identifies
	which teams operate as structurally self-contained sub-graphs (isolated DFS
	tree roots) versus those deeply embedded in larger communication patterns.
	
	%% ──────────────────────────────────────────────────────────────────────────
	\subsection{Transportation Networks and Smart Cities}
	\label{sec:transport}
	
	\paragraph{Multi-Modal Transport.}
	A city's transport system consists of multiple modal graphs (metro, bus,
	cycling, pedestrian).  HS-DFS on this collection detects which modal
	sub-networks are structurally contained within others -- e.g., whether the
	cycling network is subgraph-isomorphic to a sub-region of the road network.
	This supports:
	\begin{itemize}
		\item Intermodal journey planning (DFS tree paths give containment-based
		transfer points).
		\item Infrastructure investment prioritisation (nodes with many incoming
		meta-edges are high-leverage intervention points).
	\end{itemize}
	
	\paragraph{Smart Grid Management.}
	Power distribution networks are hierarchical graphs: transmission lines,
	distribution substations, neighbourhood grids.  HS-DFS on the collection
	of sub-grid graphs detects structural vulnerabilities (single points of
	containment whose loss would fragment the meta-graph).
	
	%% ════════════════════════════════════════════════════════════════════════════
	\section{Visualisation of the Meta-Graph}
	\label{sec:viz}
	
	\begin{figure}[H]
		\centering
		\begin{tikzpicture}[
			node/.style={circle, draw=algoblue, fill=algoblue!10, thick,
				minimum size=1.1cm, font=\small\bfseries},
			treeedge/.style={-{Stealth[length=7pt]}, thick, algoblue},
			backedge/.style={-{Stealth[length=7pt]}, thick, red!70!black, dashed},
			crossedge/.style={-{Stealth[length=7pt]}, thick, green!50!black, dotted},
			fwdedge/.style={-{Stealth[length=7pt]}, thick, orange!80!black, dash dot},
			]
			\node[node] (g1) at (0,4)   {$G_1$};
			\node[node] (g2) at (-2,2)  {$G_2$};
			\node[node] (g3) at (2,2)   {$G_3$};
			\node[node] (g4) at (-3,0)  {$G_4$};
			\node[node] (g5) at (-1,0)  {$G_5$};
			\node[node] (g6) at (1,0)   {$G_6$};
			\node[node] (g7) at (3,0)   {$G_7$};
			
			\draw[treeedge] (g1) -- (g2) node[midway,left,font=\tiny]{tree};
			\draw[treeedge] (g1) -- (g3) node[midway,right,font=\tiny]{tree};
			\draw[treeedge] (g2) -- (g4) node[midway,left,font=\tiny]{tree};
			\draw[treeedge] (g2) -- (g5) node[midway,right,font=\tiny]{tree};
			\draw[treeedge] (g3) -- (g6) node[midway,left,font=\tiny]{tree};
			\draw[treeedge] (g3) -- (g7) node[midway,right,font=\tiny]{tree};
			\draw[backedge]  (g5) to[bend right=40] (g1)
			node[midway,left,font=\tiny,red!70!black]{back};
			\draw[crossedge] (g4) to[bend right=20] (g6)
			node[midway,below,font=\tiny,green!50!black]{cross};
			\draw[fwdedge]   (g1) to[bend left=50] (g7)
			node[midway,right,font=\tiny,orange!80!black]{fwd};
			
		\end{tikzpicture}
		\caption{Schematic DFS forest on a meta-graph of seven graphs.  Tree edges
			(solid blue) form the primary containment hierarchy.  The back edge (dashed
			red, $G_5 \to G_1$) indicates circular containment.  The cross edge (dotted
			green, $G_4 \to G_6$) reveals lateral structural similarity.  The forward
			edge (dash-dot orange, $G_1 \to G_7$) encodes transitive containment.}
		\label{fig:metaforest}
	\end{figure}
	
	%% ════════════════════════════════════════════════════════════════════════════
	\section{Correctness}
	\label{sec:correctness}
	
	\begin{theorem}[Correctness of HS-DFS]
		\label{thm:correct}
		Let $\mathcal{G} = \{G_1, \ldots, G_N\}$ and let $\mathcal{M}^*$ be the
		true meta-graph (with oracle answering correctly).  HS-DFS constructs
		$\mathcal{M} = \mathcal{M}^*$ and produces a valid DFS forest
		$\mathcal{F}$ of $\mathcal{M}^*$.
	\end{theorem}
	
	\begin{proof}
		\textit{Phase~1 correctness:} The double loop enumerates all $N(N-1)$
		ordered pairs with $i \neq j$.  By \Cref{thm:subgraph_correct}, the
		Subgraph Algorithm correctly decides $G_i \hookrightarrow G_j$ for every
		pair.  Hence an edge $(G_i, G_j)$ is inserted if and only if
		$G_i \hookrightarrow G_j$, establishing $\mathcal{M} = \mathcal{M}^*$.
		
		\textit{Phase~2 correctness:} \Call{DFS-Visit}{} is the standard DFS
		procedure with colouring and timestamp assignment.  By the DFS correctness
		theorem \cite[Theorem 22.2]{cormen2009}, DFS on any directed graph produces
		a valid spanning forest, assigns monotonically increasing timestamps, and
		classifies all edges correctly.  Applying DFS to the correctly constructed
		$\mathcal{M} = \mathcal{M}^*$ therefore produces a valid DFS forest.
	\end{proof}
	
	%% ════════════════════════════════════════════════════════════════════════════
	\section{Discussion and Future Work}
	\label{sec:discussion}
	
	\paragraph{Parallelisation.}
	The dominant cost is the $N^2$ oracle invocations in Phase~1, which are
	fully independent.  A straightforward MapReduce parallelisation assigns
	$N^2 / P$ pairs to each of $P$ processors, achieving
	$\mathcal{O}(n^5 / P)$ time.  GPU-accelerated subgraph isomorphism oracles
	(e.g., using graph neural network approximators) could further reduce
	$T_{\mathcal{S}}(n)$ in practice.
	
	\paragraph{Approximate Oracles.}
	The Subgraph Algorithm \cite{epp2026subgraph} is the concrete oracle
	used throughout this paper.  Its $\mathcal{O}(n^3)$ complexity arises
	from $n$ cyclic rotation trials each requiring an $\mathcal{O}(n^2)$
	LCS computation.  For very large collections where even cubic oracle
	cost is prohibitive, one may limit the number of rotation trials to
	$\sqrt{n}$, reducing the oracle to $\mathcal{O}(n^{2.5})$ at the cost
	of a bounded false-negative rate.  Alternatively, approximate oracles
	based on Weisfeiler-Leman graph kernels or GNN embeddings yield
	probabilistic guarantees and run in $\mathcal{O}(n)$, reducing total
	complexity to $\mathcal{O}(n^3)$ with a controlled false-positive rate.
	
	\paragraph{Dynamic Collections.}
	In many applications (online KG updates, streaming genomics, live IoT
	topologies), $\mathcal{G}$ is not static.  An incremental variant of HS-DFS
	that maintains $\mathcal{M}$ under graph insertions and deletions without
	full recomputation is a natural extension.
	
	\paragraph{Weighted Meta-Graphs.}
	Assigning weights to meta-edges -- e.g., the quality or confidence of the
	subgraph embedding -- transforms the meta-graph into a weighted directed
	graph.  The DFS forest can then be replaced by a weighted DFS or
	longest-path computation to find the \emph{most significant} containment
	hierarchy.
	
	%% ════════════════════════════════════════════════════════════════════════════
	\section{Conclusion}
	\label{sec:conclusion}
	
	We have presented the \textbf{Hierarchical Subgraph-DFS (HS-DFS)} algorithm,
	a formally grounded two-phase procedure for discovering and structuring
	subgraph containment relationships within large collections of graphs.
	As its concrete oracle, HS-DFS employs the \textbf{Subgraph Algorithm}
	\cite{epp2026subgraph}: a signature-based procedure that detects pairwise
	subgraph containment in $\mathcal{O}(n^3)$ by exploiting column signature
	injectivity (\Cref{lem:injective}) and cyclic vertex re-labelling
	(\Cref{def:rotation}), entirely avoiding the factorial cost of full
	permutation search.  The combined algorithm achieves $\mathcal{O}(n^5)$
	worst-case complexity, dominated by $\mathcal{O}(n^2)$ pairwise Subgraph
	Algorithm invocations each costing $\mathcal{O}(n^3)$,
	with Phase~2 DFS running in $\mathcal{O}(n^2)$ and producing a DFS
	forest -- not achievable by BFS -- that encodes the full hierarchical, cyclic,
	and lateral structure of the collection.
	
	The scope of applicability demonstrated in \Cref{sec:applications} -- spanning
	neuroscience, genomics, logistics, industrial automation, software engineering,
	cheminformatics, semantic web technologies, social network analysis, and smart
	city infrastructure -- establishes HS-DFS as a broadly useful algorithmic
	primitive for the emerging paradigm of \emph{systems of systems}.  As
	the scale and complexity of graph-structured data continue to grow across
	all these fields, efficient algorithms for hierarchical graph analysis will
	become increasingly central to both scientific discovery and engineering
	practice.

	%% ────────────────────────────────────────────────────────────────────────────
	\subsection{Outlook}
	\label{sec:outlook}

	The present work establishes a formal foundation and a concrete algorithmic
	realisation for hierarchical subgraph analysis.  Yet it also opens a broad
	horizon of open research questions, engineering challenges, and cross-domain
	opportunities.  The following subsections outline the most promising and
	compelling directions for future investigation.

	%% · · · · · · · · · · · · · · · · · · · · · · · · · · · · · · · · · · · ·
	\subsubsection*{Reducing Oracle Complexity and Fine-Grained Hardness}

	The $\mathcal{O}(n^3)$ oracle of the Subgraph Algorithm
	\cite{epp2026subgraph} is the primary bottleneck of HS-DFS, contributing
	an $\mathcal{O}(n^5)$ overall complexity.  While this is already
	exponentially faster than exhaustive permutation search, the question of
	whether the oracle can be reduced to $o(n^3)$ remains open.  Two concrete
	sub-directions deserve attention.

	\emph{Algebraic and spectral approaches.}  Subgraph isomorphism is intimately
	related to the spectrum of the adjacency matrix: isomorphic graphs share
	identical spectra, and many containment facts are reflected in interlacing
	theorems for eigenvalues of induced subgraphs.  Recent advances in matrix
	multiplication \cite{williams2018} reduce the complexity of many
	algebraically formulated graph problems; it is conceivable that a
	reformulation of the column-signature LCS step in terms of matrix operations
	could yield an $\mathcal{O}(n^{2.37\ldots})$ oracle.  Such a result would
	immediately improve the total complexity of HS-DFS to
	$\mathcal{O}(n^{4.74\ldots})$.

	\emph{Conditional lower bounds.}  On the other hand, subgraph isomorphism
	for general graphs is NP-complete \cite{cook1971}, and even for fixed
	pattern graphs various conditional lower bounds (under, e.g., the
	Strong Exponential Time Hypothesis) have been established.  A systematic
	investigation of which graph families admit sub-cubic containment
	detection -- beyond the bounded-degree or tree-width cases already
	studied in parameterized complexity -- would clarify the fundamental
	limits of improvement and guide the design of practical heuristics.

	\emph{Heuristic pruning and early termination.}  Even without asymptotic
	improvement, substantial practical speed-ups are achievable through
	aggressive pruning: degree-sequence filtering, neighbourhood-hash
	pre-checks, and conflict-driven backtracking can reduce the number of
	LCS calls by orders of magnitude on real-world sparse graphs.  Integrating
	such heuristics rigorously into the Subgraph Algorithm while maintaining
	formal correctness guarantees is a valuable engineering task.

	%% · · · · · · · · · · · · · · · · · · · · · · · · · · · · · · · · · · · ·
	\subsubsection*{Parallelisation, GPU Acceleration, and Distributed Computation}

	The Phase~1 oracle matrix -- the $N \times N$ Boolean table recording, for
	every ordered pair $(G_i, G_j)$, whether $G_i \subseteq G_j$ -- is
	completely data-parallel: each of the $N^2$ entries can be computed
	independently.  This structure is ideally suited for massive parallelisation.

	\emph{Multi-core and SIMD parallelism.}  On a shared-memory machine with $P$
	cores, the $N^2$ oracle invocations can be distributed uniformly, reducing
	Phase~1 wall-clock time to $\mathcal{O}(n^5/P)$.  With $P = \mathcal{O}(n)$
	processors -- as available on modern many-core CPUs -- Phase~1 reduces to
	$\mathcal{O}(n^4)$.  Careful NUMA-aware load balancing and cache-oblivious
	matrix layout are necessary for efficient SIMD utilisation.

	\emph{GPU-accelerated graph kernels.}  Graphics processing units excel at
	fine-grained data-parallel workloads.  The LCS step within each oracle call
	is a standard dynamic-program\-ming recurrence that maps directly to warp-level
	parallelism.  A CUDA implementation of the Subgraph Algorithm, processing
	thousands of $(G_i, G_j)$ pairs concurrently in a single GPU kernel, is a
	straightforward and high-impact engineering contribution.  First experiments
	suggest that GPU acceleration can yield one to two orders of magnitude
	speed-up over single-threaded CPU execution for graphs of up to several
	hundred vertices.

	\emph{Distributed graph processing frameworks.}  For collections too large
	to fit in the memory of a single machine -- e.g., Knowledge Graphs with
	millions of subgraphs -- distributed frameworks such as Apache Spark
	GraphX, Pregel, or PowerGraph can be adapted.  The key insight is that
	Phase~1 reduces to a distributed matrix-vector product over the oracle
	relation, a well-studied primitive.  Phase~2 DFS can be executed on the
	aggregated meta-graph once Phase~1 completes, or, for very large
	meta-graphs, via distributed DFS variants.  Message-passing overheads
	and load-imbalance due to varying graph sizes within the collection pose
	concrete research challenges.

	\emph{Federated and privacy-preserving computation.}  In settings where
	individual graphs encode sensitive data (e.g., patient brain connectivity
	networks or proprietary molecular structures), the oracle must be evaluated
	without revealing the full graph to a central server.  Secure two-party
	computation protocols for subgraph isomorphism, possibly leveraging
	homomorphic encryption or garbled circuits, are an important direction
	at the intersection of privacy and algorithmic efficiency.

	%% · · · · · · · · · · · · · · · · · · · · · · · · · · · · · · · · · · · ·
	\subsubsection*{Approximate, Probabilistic, and Learned Oracles}

	Exact subgraph containment is a hard combinatorial predicate.  In many
	applications, a small controlled error rate is acceptable in exchange
	for dramatically reduced computation time.

	\emph{Weisfeiler-Leman graph kernels.}  The $k$-dimensional
	Weisfeiler-Leman (WL) test \cite{weisfeiler1968} provides a polynomial-time
	graph similarity measure that is strictly weaker than isomorphism but captures
	rich structural information.  A WL-based approximate oracle can be evaluated
	in $\mathcal{O}(n \log n)$ time, reducing total HS-DFS complexity to
	$\mathcal{O}(n^3 \log n)$.  The false-negative rate of such an oracle --
	i.e., pairs where $G_i \subseteq G_j$ holds but the WL kernel does not
	detect it -- depends on the graph family and can be bounded analytically for
	certain restricted classes.

	\emph{Graph Neural Network embeddings.}  Training a GNN to embed graphs
	into a Euclidean space such that $G_i \subseteq G_j$ iff
	$\mathbf{e}(G_i) \preceq \mathbf{e}(G_j)$ (in some partial-order sense) is
	an active research area.  Order-embedding models \cite{weisfeiler1968} have
	shown promise for relation learning; extending them to subgraph ordering
	would yield a learned oracle with $\mathcal{O}(n)$ inference time.
	Key challenges include training data generation, generalisation across
	graph families, and providing probabilistic guarantees on the error rate.

	\emph{Randomised algorithms with one-sided error.}  A randomised algorithm
	that correctly identifies all true subgraph containment pairs (no false
	negatives) and produces false positives with probability at most $\delta$
	would be highly valuable: the HS-DFS meta-graph may have spurious edges,
	but the DFS forest remains a valid over-approximation of the true
	containment hierarchy.  Designing such algorithms via, e.g., random
	graph fingerprinting or colour-coding \cite{cormen2009} is a natural
	extension.

	\emph{Adaptive oracle scheduling.}  Rather than committing to a single
	oracle type, an adaptive scheduler could start with a cheap approximate
	oracle and fall back to the exact Subgraph Algorithm only for pairs flagged
	as uncertain.  Bayesian active learning methods could guide which
	pairs to re-evaluate at higher cost, minimising the total computational
	budget for a given accuracy target.

	%% · · · · · · · · · · · · · · · · · · · · · · · · · · · · · · · · · · · ·
	\subsubsection*{Dynamic and Streaming Graph Collections}

	The current formulation of HS-DFS assumes a static collection
	$\mathcal{G} = \{G_1, \ldots, G_N\}$ given in advance.  Real-world
	deployments typically involve collections that evolve continuously.

	\emph{Incremental meta-graph maintenance.}  When a new graph $G_{N+1}$ is
	inserted, only the $2N$ oracle entries involving $G_{N+1}$ need to be
	recomputed (i.e., whether $G_{N+1} \subseteq G_j$ or $G_j \subseteq G_{N+1}$
	for each existing $j$).  The meta-graph can then be updated by adding a new
	node together with its incident edges.  However, the DFS forest may change
	globally even with a single edge insertion, because new containment
	relationships can introduce alternative DFS traversal orders.  Algorithms
	for \emph{incremental DFS} on directed graphs, recently developed in the
	context of online graph algorithms, provide a starting point; adapting them
	to the HS-DFS setting requires careful treatment of back-edges and
	cross-edges in the meta-graph.

	\emph{Graph deletions and expiry.}  When a graph $G_i$ is removed from the
	collection, its node and all incident edges must be deleted from
	$\mathcal{M}$.  If $G_i$ was an intermediate node in a containment chain
	$G_a \subseteq G_i \subseteq G_b$, the transitive closure of the
	meta-graph must be updated to reflect the direct relationship
	$G_a \subseteq G_b$.  Maintaining the DFS forest under deletions without
	full recomputation is a challenging open problem.

	\emph{Streaming with bounded memory.}  In resource-constrained settings
	(e.g., edge computing at IoT devices), storing the full oracle matrix is
	infeasible.  A streaming variant of HS-DFS that processes graphs in a
	single pass and maintains a compact summary -- such as a spanning subgraph
	of $\mathcal{M}$ or a compressed representation of the DFS forest -- would
	enable hierarchical analysis under strict memory constraints.

	\emph{Temporal containment.}  Many real networks carry timestamps on edges
	or evolve according to known temporal patterns (e.g., daily rhythms in brain
	activity, seasonal supply-chain patterns).  Defining \emph{temporal subgraph}
	containment -- where the embedding must respect temporal ordering -- and
	building a corresponding HS-DFS variant for temporal graph streams is an
	open problem of both theoretical and practical significance.

	%% · · · · · · · · · · · · · · · · · · · · · · · · · · · · · · · · · · · ·
	\subsubsection*{Weighted, Attributed, and Multi-Relational Graphs}

	The Subgraph Algorithm and HS-DFS as presented operate on unweighted,
	unlabelled, directed graphs.  Practical graph data is almost universally
	richer.

	\emph{Edge and vertex weights.}  Incorporating weights requires a
	generalised notion of subgraph containment: an embedding
	$\varphi: V(G_i) \to V(G_j)$ must preserve not only the edge structure
	but also weight constraints (e.g., $w_i(u,v) \leq w_j(\varphi(u),\varphi(v))$).
	The column-signature framework can be extended to encode weight ranges into
	the signature entries; LCS comparisons must then be replaced by
	order-respecting matching.

	\emph{Vertex and edge labels.}  Labelled graphs arise naturally in
	cheminformatics (atom types and bond types), semantic web (RDF predicates),
	and genomics (gene ontology annotations).  Label compatibility constraints
	add a multiplicative factor to the oracle complexity in the worst case,
	but in practice restrict the search space substantially.  The
	column-signature approach can be adapted by partitioning vertices by label
	before constructing signatures, maintaining the core algorithmic structure.

	\emph{Heterogeneous and multi-relational graphs.}  Knowledge graphs such
	as Wikidata or Freebase contain multiple distinct edge types
	(``is-a'', ``part-of'', ``located-in'', etc.).  Subgraph containment in
	such graphs -- where edge types must be matched -- is significantly harder;
	the problem is related to conjunctive query containment, which is
	$\Pi_2^P$-complete in general.  Practical algorithms exploiting the
	sparsity and regularity of real KG schemas are an important direction.

	\emph{Probabilistic and uncertain graphs.}  Graphs derived from statistical
	inference (e.g., brain connectivity from fMRI correlations, gene regulatory
	network inference) carry inherent uncertainty.  Each edge exists with some
	probability $p_{ij} \in [0,1]$.  Defining \emph{probabilistic subgraph
	containment} -- $G_i \subseteq G_j$ with probability at least $1-\epsilon$ --
	and designing HS-DFS variants that propagate uncertainty through the
	meta-graph constitute a rich open direction.

	%% · · · · · · · · · · · · · · · · · · · · · · · · · · · · · · · · · · · ·
	\subsubsection*{Hypergraphs, Simplicial Complexes, and Higher-Order Structures}

	Biological and social networks frequently exhibit \emph{higher-order
	interactions}: reactions involving three or more molecules simultaneously,
	group communications, or higher-dimensional topological features.  These
	structures are naturally represented as hypergraphs or simplicial complexes,
	neither of which is captured by the ordinary graph model.

	\emph{Hypergraph subgraph containment.}  A hyperedge connects an arbitrary
	subset of vertices; a hypergraph $H_i$ is a sub-hypergraph of $H_j$ if
	there exists an injective map on vertices that sends each hyperedge of
	$H_i$ to a subset of a hyperedge of $H_j$.  The column-signature approach
	does not directly generalise, since adjacency matrices are replaced by
	incidence matrices with variable row sizes.  A tensorial encoding of
	hyperedge memberships may provide a path forward.

	\emph{Simplicial complex homology.}  Persistent homology -- tracking the
	birth and death of topological features (connected components, loops,
	voids) across a filtration -- provides a summary of the hierarchical
	structure of a simplicial complex.  Relating the HS-DFS forest to the
	persistent homology barcode of the collection is a deep and speculative
	but potentially very fruitful direction: the DFS forest encodes
	combinatorial containment, while barcodes encode topological containment.
	Understanding when and how these two notions align would illuminate the
	structural basis of many hierarchical phenomena in nature.

	%% · · · · · · · · · · · · · · · · · · · · · · · · · · · · · · · · · · · ·
	\subsubsection*{Integration with Graph Neural Networks and Deep Learning}

	Graph Neural Networks (GNNs) have revolutionised machine learning on
	graph-structured data.  Their integration with HS-DFS opens several
	mutually beneficial directions.

	\emph{HS-DFS as a supervision signal.}  The containment hierarchy produced
	by HS-DFS can serve as structured supervision for training GNNs to learn
	\emph{hierarchical graph representations}: representations in which
	$\mathbf{e}(G_i)$ is ``below'' $\mathbf{e}(G_j)$ in embedding space
	whenever $G_i \subseteq G_j$.  Such representations would enable fast
	approximate lookups in large graph databases without recomputing the oracle.
	The DFS tree structure provides a natural curriculum: simpler containment
	relationships (few vertices) can be used for initial training, with
	progressively deeper relationships added as training proceeds.

	\emph{GNN expressiveness and the WL hierarchy.}  The expressive power of
	standard GNNs is bounded by the 1-dimensional Weisfeiler-Leman test
	\cite{weisfeiler1968}.  Since the Subgraph Algorithm uses column signatures
	derived from the adjacency matrix -- which go beyond 1-WL in their
	sensitivity to global structure via cyclic rotations -- it offers a
	concrete candidate for a ``Subgraph-WL'' test that may lie strictly between
	1-WL and 2-WL in discriminative power.  Formalising this relationship and
	designing GNN architectures that match it would contribute to the ongoing
	programme of characterising GNN expressiveness.

	\emph{Meta-graph embeddings for graph retrieval.}  Given a query graph $Q$
	and a large database $\mathcal{G}$, finding all $G_i \in \mathcal{G}$ with
	$Q \subseteq G_i$ is the \emph{supergraph retrieval} problem, central to
	drug discovery and code search.  The meta-graph $\mathcal{M}$ provides an
	index that dramatically prunes the search: once $Q$ is located (or
	approximately located) in $\mathcal{M}$, only ancestors in the DFS forest
	need to be checked.  Training a GNN to predict query position in
	$\mathcal{M}$ without constructing $\mathcal{M}$ explicitly is a promising
	direction for scalable graph retrieval.

	\emph{Transformer architectures on meta-graphs.}  Attention-based graph
	transformers operating on $\mathcal{M}$ as the input graph -- with nodes
	as graphs and edges as containment relations -- can learn complex
	meta-level patterns: e.g., which structural motifs are universally present
	across a graph collection, or which subgraphs tend to co-occur.  This
	``graph of graphs'' view enables a new class of second-order graph learning
	problems that HS-DFS is uniquely positioned to support.

	%% · · · · · · · · · · · · · · · · · · · · · · · · · · · · · · · · · · · ·
	\subsubsection*{Quantum Algorithms for Subgraph Problems}

	Quantum computing offers the prospect of qualitatively different complexity
	trade-offs for hard combinatorial problems.

	\emph{Grover-based oracle speedup.}  Grover's algorithm provides a
	quadratic speedup for unstructured search.  Applied to the permutation
	search underlying naive subgraph isomorphism, it reduces the search space
	from $\mathcal{O}(n!)$ to $\mathcal{O}(\sqrt{n!})$ -- still
	super-exponential, but a meaningful improvement.  Whether the
	column-signature structure of the Subgraph Algorithm admits a quantum
	analogue that achieves better than $\mathcal{O}(n^{1.5})$ oracle
	complexity is an intriguing open question.

	\emph{Quantum graph isomorphism.}  Recent work on quantum query complexity
	has established non-trivial bounds for graph isomorphism; extending these
	results to the subgraph setting is largely unexplored.  The connection
	between the cyclic rotation invariance exploited in the Subgraph Algorithm
	and quantum Fourier transforms over cyclic groups is a structural hint
	that quantum methods might be particularly natural here.

	\emph{Variational quantum eigensolver approaches.}  For small-to-medium
	graphs, variational quantum algorithms (VQE, QAOA) can be used to encode
	the subgraph isomorphism problem as a quadratic unconstrained binary
	optimisation (QUBO) instance and approximately solve it on near-term
	quantum hardware.  Although current hardware is too noisy for a practical
	advantage, this direction will become increasingly relevant as quantum
	processors scale.

	%% · · · · · · · · · · · · · · · · · · · · · · · · · · · · · · · · · · · ·
	\subsubsection*{Parameterized Complexity and Structural Graph Theory}

	The classical intractability of subgraph isomorphism is alleviated by
	\emph{parameterized complexity theory}, which identifies structural
	parameters of the input that, when bounded, admit efficient algorithms.

	\emph{Tree-width and other width parameters.}  For pattern graphs of
	bounded tree-width $tw$, subgraph isomorphism can be solved in
	$f(tw) \cdot n^{\mathcal{O}(1)}$ time via Courcelle's theorem (a result
	in meta-algorithmic form) or more concretely via tree-decomposition
	dynamic programming.  It is natural to ask whether the Subgraph Algorithm
	can be adapted to exploit bounded tree-width of either the pattern or the
	host graph, reducing oracle complexity below $\mathcal{O}(n^3)$ in this
	important special case.

	\emph{Graph minors and excluded subgraph families.}  Robertson-Seymour
	theory guarantees that any minor-closed graph family is characterised by a
	finite set of excluded minors.  For graph collections arising from
	specific applications (e.g., planar graphs in geographic networks, outerplanar
	graphs in certain chemical bond structures), the excluded-minor characterisation
	may enable polynomial-time exact subgraph testing far faster than the general
	$\mathcal{O}(n^3)$ oracle.

	\emph{Colour coding and randomized FPT algorithms.}  For pattern graphs of
	bounded size $k$, colour-coding \cite{cormen2009} achieves
	$\mathcal{O}(2^k \cdot n^{\mathcal{O}(1)})$ running time for subgraph
	detection.  When the collection $\mathcal{G}$ consists of small graphs
	(e.g., $k = \mathcal{O}(\log n)$), this significantly outperforms the
	cubic oracle.  Integrating color-coding into HS-DFS for the regime of
	small graphs is a concrete improvement opportunity.

	%% · · · · · · · · · · · · · · · · · · · · · · · · · · · · · · · · · · · ·
	\subsubsection*{Multi-Layer Networks and Interdependent Systems}

	Modern infrastructure, social systems, and biological organisms are not
	captured by a single network: they consist of \emph{multiple layers} of
	interactions that are coupled to one another.

	\emph{Multi-layer subgraph containment.}  A multi-layer graph
	$\mathbf{G} = (G^{(1)}, \ldots, G^{(L)})$ is a tuple of graphs,
	one per layer, sharing the same vertex set.  A natural generalisation
	of subgraph containment requires an embedding that is valid simultaneously
	in all layers.  HS-DFS on collections of multi-layer graphs would produce
	a meta-graph whose edges certify \emph{joint} structural containment --
	a far stronger relation than containment in any single layer.

	\emph{Coupled DFS forest.}  In the DFS phase, the DFS forest of the
	multi-layer meta-graph may couple containment hierarchies that are
	independent within each layer.  Understanding the resulting structure --
	when does a node occupy the same position across all layer hierarchies,
	and when does layer coupling create new hierarchy patterns? -- is both
	theoretically interesting and practically relevant for, e.g., analysing
	multiplex social networks or multi-omics biological data.

	\emph{Application to critical infrastructure.}  Power grids, communication
	networks, and transportation systems are tightly interdependent.  Their
	joint failure modes are governed by \emph{cascading failures} across layers.
	HS-DFS applied to a collection of multi-layer subgraphs representing possible
	failure scenarios could identify which scenarios are structurally contained
	in others, enabling a partial order over risk scenarios and guiding
	prioritised mitigation strategies.

	%% · · · · · · · · · · · · · · · · · · · · · · · · · · · · · · · · · · · ·
	\subsubsection*{Meta-Graph Visualisation and Interactive Exploration}

	\Cref{sec:visualisation} introduced initial techniques for rendering the
	meta-graph.  As $N$ and the complexity of $\mathcal{M}$ grow, scalable and
	meaningful visualisation becomes a significant challenge in its own right.

	\emph{Hierarchical layout algorithms.}  The DFS forest is a natural
	scaffold for a hierarchical graph layout: DFS roots form the topmost level,
	DFS children descend, and cross-edges and back-edges are rendered as
	additional arcs.  Adapting the Sugiyama layered graph drawing framework
	to exploit the DFS forest structure -- in particular, to minimise
	edge crossings while preserving the forest's tree sekeleton -- is a
	concrete visualisation research goal.

	\emph{Interactive drill-down and filtering.}  A visual analytics tool
	built on top of HS-DFS should allow a user to click on any node in the
	meta-graph to inspect the corresponding graph, to filter the view to
	a subgraph of $\mathcal{M}$ induced by a vertex subset of interest, and
	to highlight the shortest containment chain between two selected graphs.
	Integration with existing graph visualisation libraries (Gephi, Cytoscape,
	D3.js) would facilitate rapid prototyping.

	\emph{Summary statistics and meta-graph dashboards.}  For very large
	collections ($N > 10^4$), interactive node-link rendering is impractical.
	Statistical summaries -- degree distribution of $\mathcal{M}$, DFS forest
	depth distribution, fraction of DAG vs.\ tree edges -- can be displayed
	as dashboards that convey the global shape of the containment hierarchy
	without rendering every node.  Anomaly detection on these statistics
	could flag unexpected structural changes in a dynamically evolving
	collection.

	\emph{Augmented and virtual reality.}  For collections arising from
	high-dimensional scientific data (e.g., full connectomes with thousands
	of brain regions, or molecular databases with tens of thousands of
	compounds), immersive 3D exploration in augmented or virtual reality
	environments may offer cognitive advantages over flat 2D representations.
	Mapping the DFS forest depth to a physical $z$-coordinate and using
	spatial audio to encode edge density could exploit human depth perception
	for rapid hierarchical exploration.

	%% · · · · · · · · · · · · · · · · · · · · · · · · · · · · · · · · · · · ·
	\subsubsection*{Benchmarking, Open Datasets, and Reproducibility}

	The empirical evaluation of HS-DFS and competing approaches requires
	standardised benchmarks that do not yet exist in a unified form.

	\emph{Synthetic benchmark suites.}  A systematic benchmark suite should
	cover the full range of structural parameters: varying $N$ (collection
	size), $n$ (individual graph size), graph density (sparse to dense),
	degree of nesting (fraction of true containment pairs), and graph
	regularity (ranging from random Erd\H{o}s-R\'enyi graphs to structured
	lattices and trees).  All benchmarks should be generated by reproducible
	random models with fixed seeds and published under open licences.

	\emph{Real-world dataset curation.}  Each of the application domains
	discussed in \Cref{sec:applications} provides potential real-world
	datasets: functional connectome archives (HCP, UK Biobank), metabolic
	network databases (KEGG \cite{kegg}, BiGG), software dependency graphs
	(Maven Central, npm), molecular databases (ChEMBL, PubChem), and
	knowledge graph snapshots (Wikidata dumps).  Curating graph collections
	from these sources, annotating ground-truth containment relationships
	where available, and publishing them as a community benchmark would
	greatly accelerate research on hierarchical graph algorithms.

	\emph{Reproducibility and open-source tooling.}  The Python implementation
	of HS-DFS accompanying this paper \cite{epp2026subgraph} provides a
	reference baseline.  Future work should encompass:
	a comprehensive test suite with property-based testing;
	continuous integration pipelines that enforce correctness on every commit;
	containerised environments (Docker, Singularity) for noise-free timing
	experiments;
	and publication of all experimental artefacts (code, data, logs) under
	open licences.  Adherence to the FAIR principles (Findable, Accessible,
	Interoperable, Reusable) for research software is essential for meaningful
	comparative evaluation.

	%% · · · · · · · · · · · · · · · · · · · · · · · · · · · · · · · · · · · ·
	\subsubsection*{Domain-Specific Optimisations and Deployment}

	Beyond algorithmic generalisation, each of the application domains from
	\Cref{sec:applications} offers specific opportunities to tailor and deploy
	HS-DFS in practice.

	\emph{Neuroscience and connectomics.}  The brain's connectome is sparse,
	hierarchically organised, and exhibits strong bilateral symmetry.  Exploiting
	these priors -- for instance, restricting the oracle to
	connectivity-preserving embeddings or enforcing hemisphere symmetry
	constraints -- can reduce oracle complexity by large constant factors.
	A dedicated HS-DFS module integrated into existing connectome analysis
	platforms (e.g., the Human Connectome Project pipeline or Nipype) would
	enable large-scale comparative connectomics.

	\emph{Drug discovery and cheminformatics.}  Molecular graphs contain at
	most a few hundred atoms; $n$ is small but $N$ (number of compounds in a
	screen) can reach $10^9$ in virtual screening campaigns.  An HS-DFS-based
	scaffold-tree decomposition, computing which molecular fragments are
	subgraphs of which larger molecules, would provide a principled alternative
	to existing heuristic scaffold-tree tools (e.g., Scaffold Hunter, RDKit
	Murcko decomposition) with formal correctness guarantees.

	\emph{Software engineering and program analysis.}  Control-flow and
	call-graph containment between software modules provides a principled
	basis for \emph{library recommendation}: if a project's dependency graph
	is a subgraph of that of a well-tested reference project, specific
	library versions are formally compatible.  Integrating HS-DFS into
	package managers (Cargo, pip, npm) as an optional compatibility checker --
	invoked at install time to verify structural compatibility -- is a
	concrete deployment scenario.

	\emph{Knowledge graph completion and alignment.}  In the semantic web,
	two ontologies $O_1$ and $O_2$ are \emph{aligned} if there exists a
	containment embedding between their concept graphs.  HS-DFS provides
	not only a binary alignment verdict but the full DFS forest, which
	encodes the alignment structure at all levels of granularity.  Integrating
	HS-DFS with existing ontology alignment tools (LogMap, AgreementMaker Light)
	as a post-processing step to verify and structure proposed alignments is
	a concrete short-term engineering goal.

	%% · · · · · · · · · · · · · · · · · · · · · · · · · · · · · · · · · · · ·
	\subsubsection*{Theoretical Foundations: Algebra, Order Theory, and Logic}

	Beyond algorithmic questions, HS-DFS invites a range of deeper
	theoretical investigations.

	\emph{The poset of graphs under subgraph containment.}  The subgraph
	containment relation (modulo isomorphism) defines a partial order on
	the set of all finite graphs.  The meta-graph $\mathcal{M}$ is a
	finite induced sub-poset of this infinite structure.  Understanding the
	order-theoretic properties of this sub-poset -- its width, its height,
	the existence of antichains of prescribed size -- is a problem in
	combinatorics with direct implications for the complexity of HS-DFS
	(e.g., antichains imply no containment edges, making Phase~1 wasteful).
	Ramsey-theoretic arguments can be used to establish bounds on antichain
	size for graph collections with prescribed density.

	\emph{Modal and description logics.}  Subgraph containment can be expressed
	in monadic second-order logic (MSO): $G_i \subseteq G_j$ iff there exists
	an injective function on vertices satisfying the edge-preservation axiom.
	By Courcelle's theorem, MSO-definable properties are decidable in
	linear time on graphs of bounded tree-width.  Exploring which higher-level
	properties of the meta-graph (e.g., the number of strongly connected
	components in $\mathcal{M}$) are MSO-definable and exploiting this for
	declarative query processing is a direction at the intersection of logic
	and algorithms.

	\emph{Category-theoretic formulation.}  Graph homomorphisms and subgraph
	embeddings are natural morphisms in the category of graphs.  The
	meta-graph $\mathcal{M}$ is then a representation of part of this
	category's morphism structure.  A categorical framework for HS-DFS --
	expressing Phase~1 as a functor and the DFS forest as a skeleton of
	the resulting category -- could yield deep structural insights and
	connections to topos theory, sheaves over graph categories, and
	compositional approaches to graph rewriting.

	%% · · · · · · · · · · · · · · · · · · · · · · · · · · · · · · · · · · · ·
	\subsubsection*{Summary of the Research Horizon}

	Taken together, the directions outlined above span the full spectrum from
	immediate engineering improvements -- parallel implementations, benchmark
	datasets, domain-specific integrations -- to deep theoretical questions
	at the frontier of algorithms, complexity theory, quantum computing, and
	mathematical foundations.  The unifying theme is that the hierarchical
	containment structure captured by HS-DFS is not merely an algorithmic
	artefact but a fundamental aspect of how complex systems are organised:
	smaller functional units are embedded in larger ones, and understanding
	this embeddability at scale is essential for science and engineering alike.
	We therefore anticipate that HS-DFS, and the broader research programme
	it represents, will stimulate sustained work across multiple communities
	over the coming years.

	%% ════════════════════════════════════════════════════════════════════════════
	\begin{thebibliography}{99}
		
		\bibitem{epp2026subgraph}
		S.~Epp.
		Der Subgraph Algorithmus. Available at \url{https://github.com/hjstephan86/subgraph}, 2026.
		
		\bibitem{cormen2009}
		T.~H. Cormen, C.~E. Leiserson, R.~L. Rivest, and C.~Stein.
		\textit{Introduction to Algorithms}, 3rd~ed.
		MIT Press, Cambridge, MA, 2009.
		
		\bibitem{cook1971}
		S.~A. Cook.
		The complexity of theorem-proving procedures.
		In \textit{Proc. 3rd ACM STOC}, pp.~151--158, 1971.
		
		\bibitem{williams2018}
		V.~V. Williams.
		On some fine-grained questions in algorithms and complexity.
		In \textit{Proc. ICM}, vol.~3, pp.~3431--3472, 2018.
		
		\bibitem{ullmann1976}
		J.~R. Ullmann.
		An algorithm for subgraph isomorphism.
		\textit{J. ACM}, 23(1):31--42, 1976.
		
		\bibitem{vf2}
		L.~P. Cordella, P.~Foggia, C.~Sansone, and M.~Vento.
		A (sub)graph isomorphism algorithm for matching large graphs.
		\textit{IEEE TPAMI}, 26(10):1367--1372, 2004.
		
		\bibitem{weisfeiler1968}
		B.~Weisfeiler and A.~Leman.
		The reduction of a graph to canonical form and the algebra which appears therein.
		\textit{NTI, Series 2}, 9:12--16, 1968.
		
		\bibitem{sporns2011}
		O.~Sporns.
		\textit{Networks of the Brain}.
		MIT Press, 2011.
		
		\bibitem{kegg}
		M.~Kanehisa and S.~Goto.
		KEGG: Kyoto Encyclopedia of Genes and Genomes.
		\textit{Nucleic Acids Research}, 28(1):27--30, 2000.
		
		\bibitem{autosar}
		AUTOSAR Consortium.
		\textit{AUTOSAR Layered Software Architecture}, Release 4.4, 2019.
		
		\bibitem{iit}
		G.~Tononi.
		Consciousness as integrated information: a provisional manifesto.
		\textit{Biol. Bull.}, 215(3):216--242, 2008.
		
		\bibitem{epp2026brain}
		S.~Epp.
		Die epistemische Wolke, der unsichtbare Geist des Lebens und die
		Rechtsdrehung kortikaler Informationsverarbeitung.
		Available at \url{https://github.com/hjstephan86/science/tree/main/science/brn}, 2026.
		
		\bibitem{hodgkin1952}
		A.~L. Hodgkin and A.~F. Huxley.
		A quantitative description of membrane current and its application to
		conduction and excitation in nerve.
		\textit{Journal of Physiology}, 117(4):500--544, 1952.
		
		\bibitem{delcastillo1954}
		J.~Del Castillo and B.~Katz.
		Quantal components of the end-plate potential.
		\textit{Journal of Physiology}, 124(3):560--573, 1954.
		
		\bibitem{goldmanrakic1995}
		P.~S. Goldman-Rakic.
		Cellular basis of working memory.
		\textit{Neuron}, 14(3):477--485, 1995.
		
		\bibitem{radley2004}
		J.~J. Radley, H.~M. Sisti, J.~Hao, A.~B. Rocher, T.~McCall, P.~R. Hof,
		B.~S. McEwen, and J.~H. Morrison.
		Chronic behavioral stress induces apical dendritic reorganization in
		pyramidal neurons of the medial prefrontal cortex.
		\textit{Cerebral Cortex}, 14(7):711--717, 2004.
		
		\bibitem{hopfield1982}
		J.~J. Hopfield.
		Neural networks and physical systems with emergent collective
		computational abilities.
		\textit{Proceedings of the National Academy of Sciences},
		79(8):2554--2558, 1982.
		
	\end{thebibliography}
	
\end{document}
